\documentclass[10pt,a4paper]{article}
\usepackage[utf8]{inputenc}
\usepackage[spanish]{babel}
\usepackage{amsmath,amssymb,amsfonts}
\usepackage{geometry}
\usepackage{xcolor}
\usepackage{enumitem}
\usepackage{tcolorbox}
\usepackage{hyperref}

\geometry{top=2cm,left=2cm,right=2cm,bottom=2cm}
\setlist{leftmargin=*}
\pagestyle{plain}
\setlength{\parindent}{0pt}
\setlength{\parskip}{4pt}

\definecolor{azuloscuro}{RGB}{0,51,102}
\definecolor{azulclaro}{RGB}{51,153,255}
\definecolor{rojoclaro}{RGB}{255,102,102}
\definecolor{verdeclaro}{RGB}{102,204,102}

\hypersetup{
    colorlinks=true,
    linkcolor=azuloscuro,
    urlcolor=azuloscuro,
    citecolor=azuloscuro
}

\title{\textbf{\LARGE GUÍA COMPLETA DE COMUNICACIONES DIGITALES}\\
\Large Teoría Explicada con Detalle\\
\large TAF STAR UEC B - IMT Atlantique 2024-2025}
\author{}
\date{}

\begin{document}

\maketitle

\tableofcontents
\newpage

\section{Fundamentos}

\subsection{¿Qué son las comunicaciones digitales?}

Las comunicaciones digitales permiten transmitir información discreta (bits: 0 y 1) a través de un medio físico continuo (ondas electromagnéticas, señales eléctricas, luz en fibra óptica).

\textbf{Problema fundamental:} Necesitamos enviar datos digitales usando señales analógicas, minimizando errores y maximizando la velocidad de transmisión.

\subsection{Cadena de comunicación completa}

\textbf{Emisor:}
\begin{enumerate}
\item \textbf{Codificación de fuente:} Compresión (eliminar redundancia)
\item \textbf{Codificación de canal:} Añadir redundancia controlada para detectar/corregir errores
\item \textbf{Modulación:} Convertir bits en señales analógicas que se pueden transmitir
\item \textbf{Filtrado:} Adaptar la señal al canal (limitar ancho de banda)
\end{enumerate}

\textbf{Canal físico:}
\begin{itemize}
\item Añade \textbf{ruido aleatorio} (AWGN - Additive White Gaussian Noise)
\item Introduce \textbf{interferencias} de otras señales
\item Causa \textbf{distorsión} (multipaso, eco, reflexiones)
\item Produce \textbf{atenuación} (pérdida de potencia)
\end{itemize}

\textbf{Receptor:}
\begin{enumerate}
\item \textbf{Filtro adaptado:} Maximiza la relación señal/ruido (SNR)
\item \textbf{Muestreo:} Convierte señal continua en muestras discretas
\item \textbf{Decisión:} Determina qué símbolo se transmitió (estima bits)
\item \textbf{Decodificación:} Detecta y corrige errores usando redundancia
\end{enumerate}

\textbf{Objetivo global:} Transmitir la mayor cantidad de bits por segundo (débito) con la menor tasa de error posible, usando el mínimo ancho de banda y potencia.

\subsection{Transformada de Fourier - Herramienta fundamental}

La Transformada de Fourier (FT) es esencial porque permite analizar señales en dos dominios complementarios:

\textbf{Dominio temporal $h(t)$:} Cómo varía la señal en el tiempo

\textbf{Dominio frecuencial $H(\nu)$:} Qué frecuencias contiene la señal y con qué amplitud

$$H(\nu) = \int_{\mathbb{R}} h(t)e^{-i2\pi\nu t}dt, \quad h(t) = \int_{\mathbb{R}} H(\nu)e^{i2\pi\nu t}d\nu$$

\textbf{¿Por qué es importante?}
\begin{itemize}
\item El ancho de banda del canal limita qué frecuencias pueden pasar
\item El filtrado es multiplicación en frecuencia (más simple que convolución)
\item Nos permite diseñar señales que usen eficientemente el espectro disponible
\end{itemize}

\textbf{Propiedades clave:}
\begin{align*}
\text{FT}[h * g] &= H(\nu)G(\nu) \quad \text{(convolución → producto)}\\
\text{FT}[h \cdot g] &= H(\nu) * G(\nu) \quad \text{(producto → convolución)}\\
\text{FT}[h(t-a)] &= H(\nu)e^{-i2\pi\nu a} \quad \text{(desplazamiento temporal)}\\
\text{FT}[h(t)e^{i2\pi\nu_0 t}] &= H(\nu-\nu_0) \quad \text{(modulación → desplazamiento frecuencial)}
\end{align*}

\textbf{Teorema de Parseval (conservación de energía):}
$$E_s = \int_{\mathbb{R}} |s(t)|^2dt = \int_{\mathbb{R}} |S(\nu)|^2d\nu$$

La energía total es la misma en tiempo y frecuencia.

\subsection{Identidades útiles}

\begin{align*}
\cos^2(a) &= \frac{1 + \cos(2a)}{2}, \quad \sin^2(a) = \frac{1 - \cos(2a)}{2}\\
\cos(a)\cos(b) &= \frac{1}{2}[\cos(a-b) + \cos(a+b)]\\
\sin(a)\sin(b) &= \frac{1}{2}[\cos(a-b) - \cos(a+b)]\\
\text{sinc}(x) &= \frac{\sin(\pi x)}{\pi x}
\end{align*}

\section{Modulaciones Lineales - Convertir bits en señales}

\subsection{¿Por qué modular?}

\textbf{Problema:} Los bits son símbolos discretos (0 y 1), pero los canales físicos transmiten señales continuas.

\textbf{Solución:} La modulación asigna a cada grupo de bits una forma de onda específica que se puede transmitir.

\textbf{Dos enfoques principales:}
\begin{enumerate}
\item \textbf{Banda base:} Transmitir directamente (cables, Ethernet)
\item \textbf{Paso de banda:} Modular sobre una portadora de alta frecuencia (radio, WiFi, 4G/5G)
\end{enumerate}

\subsection{Señal transmitida}

\textbf{Banda base (sin portadora):}
$$s(t) = \sum_{k} a_k h(t-kT_s)$$

donde:
\begin{itemize}
\item $a_k$: símbolo de información en el instante $k$ (puede tomar $M$ valores diferentes)
\item $h(t)$: pulso de forma (define cómo se "dibuja" cada símbolo en el tiempo)
\item $T_s$: período de símbolo (tiempo entre símbolos consecutivos)
\end{itemize}

\textbf{Interpretación:} Cada símbolo $a_k$ "moldea" un pulso $h(t)$ que se envía en su intervalo de tiempo.

\textbf{Con portadora (forma real) - Portadora única:}
$$s(t) = \sum_{k} a_k h(t-kT_s)\cos(2\pi\nu_0 t) - b_k h(t-kT_s)\sin(2\pi\nu_0 t)$$

\textbf{¿Por qué dos componentes?} Usamos dos dimensiones ortogonales:
\begin{itemize}
\item Componente \textbf{en fase (I)}: multiplicada por $\cos(2\pi\nu_0 t)$
\item Componente \textbf{en cuadratura (Q)}: multiplicada por $-\sin(2\pi\nu_0 t)$
\end{itemize}

Esto permite transmitir el doble de información sin aumentar el ancho de banda.

\textbf{Notación compleja (más compacta):}
$$s_c(t) = \sum_{k} d_k h(t-kT_s), \quad \text{con } d_k = a_k + ib_k$$

La señal real se obtiene como: $s(t) = \Re[s_c(t)e^{i(2\pi\nu_0 t + \theta_0)}]$

\textbf{Ventajas de la portadora:}
\begin{itemize}
\item Adaptación a las características del canal (frecuencias que se propagan mejor)
\item Multiplexación (varios usuarios en diferentes frecuencias)
\item Antenas eficientes (tamaño relacionado con longitud de onda)
\end{itemize}

\subsection{Parámetros fundamentales - Métricas de eficiencia}

\begin{tcolorbox}[colback=azulclaro!10,colframe=azuloscuro,title=Definiciones clave]
\textbf{Rapidez de modulación (Symbol rate):} $D_s = \frac{1}{T_s}$ [bauds o símbolos/s]

Es la velocidad a la que enviamos símbolos. Determina el ancho de banda necesario.

\textbf{Débito binario (Bit rate):} $D_b = D_s \cdot \log_2(M) = \frac{k}{T_s}$ [bits/s]

Es la velocidad real de transmisión de información, donde $M = 2^k$ es el número de símbolos diferentes.

\textbf{Relación fundamental:} $D_b = D_s \cdot \log_2(M)$

\textbf{Ejemplo:} Si $D_s = 1000$ baudios y usamos $M=16$ símbolos (16-QAM):
$$D_b = 1000 \cdot \log_2(16) = 1000 \cdot 4 = 4000 \text{ bits/s}$$

\textbf{Ancho de banda ocupado:}
\begin{itemize}
\item Portadora única (QAM/PSK): $B_{occ} = \frac{1+\beta}{T_s} = (1+\beta)D_s$
\item PAM/ASK: $B_{occ} = \frac{2(1+\beta)}{T_s}$ (ocupa el doble)
\end{itemize}

donde $\beta$ es el factor de roll-off del filtro (exceso de banda).

\textbf{Eficiencia espectral:} $\eta = \frac{D_b}{B_{occ}}$ [bits/s/Hz]

Mide cuántos bits transmitimos por cada Hz de ancho de banda usado.

\begin{itemize}
\item Portadora única: $\eta = \frac{\log_2(M)}{1+\beta}$
\item PAM: $\eta = \frac{2\log_2(M)}{1+\beta}$
\end{itemize}

\textbf{Dilema fundamental:}
\begin{itemize}
\item Mayor $M$ → Mayor $\eta$ (más bits por segundo con mismo ancho de banda)
\item Pero Mayor $M$ → Símbolos más cercanos → Mayor probabilidad de error
\item Solución: Necesitamos mayor SNR (más potencia de transmisión)
\end{itemize}
\end{tcolorbox}

\subsection{Alfabetos de modulación - Constelaciones}

Una constelación es el conjunto de todos los símbolos posibles representados como puntos en el plano complejo (o en la recta real para PAM).

\textbf{Criterios de diseño:}
\begin{enumerate}
\item \textbf{Maximizar distancia entre símbolos} (menor tasa de error)
\item \textbf{Minimizar energía promedio} (menor consumo de potencia)
\item \textbf{Facilidad de implementación} (hardware/software)
\end{enumerate}

\textbf{M-PAM (Pulse Amplitude Modulation) - Unidimensional:}
$$\mathcal{A} = \{-(M-1)\sqrt{\varepsilon}, -(M-3)\sqrt{\varepsilon}, \ldots, -\sqrt{\varepsilon}, \sqrt{\varepsilon}, \ldots, (M-1)\sqrt{\varepsilon}\}$$

\begin{itemize}
\item Símbolos equiespaciados en el eje real (solo amplitud varía)
\item $E_{avg} = \frac{M^2-1}{3}\varepsilon$, \quad $d_{min} = 2\sqrt{\varepsilon}$
\item $\rho = \frac{d_{min}^2}{E_{avg}} = \frac{6}{M^2-1}$ (eficiencia de potencia)
\item \textbf{Uso:} Transmisión banda base (Ethernet, DSL)
\item \textbf{M-ASK:} Versión con portadora de M-PAM
\end{itemize}

\textbf{Ejemplo 4-PAM:} Símbolos en $\{-3, -1, +1, +3\}$ (normalizados)

\textbf{M-PSK (Phase Shift Keying) - Amplitud constante:}
$$\mathcal{A} = \{e^{i2\pi k/M}\sqrt{E_s} : k = 0,1,\ldots,M-1\}$$

\begin{itemize}
\item Símbolos distribuidos uniformemente en círculo de radio $\sqrt{E_s}$
\item Solo la fase varía, amplitud constante
\item $E_{avg} = E_s$, \quad $d_{min} = 2\sqrt{E_s}\sin(\pi/M)$
\item $\rho = 4\sin^2(\pi/M)$, \quad PAPR = 1 (potencia constante)
\item \textbf{Ventaja clave:} PAPR=1 → ideal para amplificadores no lineales (satélite)
\item \textbf{Desventaja:} Para $M$ grande, $d_{min}$ decrece rápidamente
\item \textbf{Uso:} Satélite (QPSK, 8-PSK), comunicaciones móviles
\end{itemize}

\textbf{Casos especiales:}
\begin{itemize}
\item BPSK ($M=2$): símbolos en $\{+\sqrt{E_s}, -\sqrt{E_s}\}$ → muy robusto
\item QPSK ($M=4$): símbolos en $\{\pm 1 \pm i\}\frac{\sqrt{E_s}}{\sqrt{2}}$ → muy usado
\end{itemize}

\textbf{M-QAM (Quadrature Amplitude Modulation) - Bidimensional:}
$$d_k = (a_k + ib_k), \quad a_k, b_k \in \{\pm 1, \pm 3, \ldots, \pm(\sqrt{M}-1)\}\sqrt{\varepsilon}$$

\begin{itemize}
\item Constelación rectangular: varía amplitud Y fase
\item $E_{avg} = \frac{2(M-1)}{3}\varepsilon$, \quad $d_{min} = 2\sqrt{\varepsilon}$
\item $\rho = \frac{6}{M-1}$, \quad $\text{PAPR} = \frac{3(\sqrt{M}-1)}{\sqrt{M}+1}$ (variable)
\item \textbf{Ventaja:} Mayor eficiencia energética que PSK para $M$ grande
\item \textbf{Razón:} Símbolos mejor distribuidos → mayor $d_{min}$ con misma $E_{avg}$
\item \textbf{Desventaja:} PAPR > 1 (necesita amplificadores lineales)
\item \textbf{Uso:} WiFi (64-QAM, 256-QAM), 4G/5G (hasta 256-QAM), cable módem
\end{itemize}

\textbf{Comparación PSK vs QAM:}
\begin{itemize}
\item Para $M$ pequeño ($M \leq 4$): PSK ≈ QAM en rendimiento
\item Para $M$ grande ($M \geq 16$): QAM superior (ganancia $\approx 3-4$ dB)
\item PSK mantiene PAPR=1 → mejor para amplificadores baratos/no lineales
\item QAM requiere amplificador lineal pero es más eficiente espectralmente
\end{itemize}

\subsection{Gray Mapping - Minimizar errores de bit}

\begin{tcolorbox}[colback=azulclaro!10,colframe=azuloscuro,title=Principio]
El \textbf{código Gray} asigna códigos binarios a símbolos de forma que símbolos adyacentes (vecinos más cercanos en la constelación) difieren en \textbf{exactamente 1 bit}.

\textbf{¿Por qué es importante?}

Cuando hay ruido, es más probable confundir un símbolo con sus vecinos cercanos. Con Gray mapping:
\begin{itemize}
\item Error de símbolo al vecino más cercano → Solo 1 bit erróneo
\item Sin Gray mapping → Pueden ser múltiples bits erróneos
\end{itemize}

\textbf{Consecuencia:} Con Gray mapping y alto SNR:
$$P_{be} \approx \frac{P_{se}}{\log_2(M)}$$

La probabilidad de error de bit es aproximadamente $\log_2(M)$ veces menor que la de símbolo.

\textbf{Ejemplo 8-PSK:}
\begin{center}
\begin{tabular}{|c|c|c|}
\hline
Fase & Símbolo & Gray \\
\hline
$0°$ & $e^{i\cdot 0}$ & 000 \\
$45°$ & $e^{i\pi/4}$ & 001 \\
$90°$ & $e^{i\pi/2}$ & 011 \\
$135°$ & $e^{i3\pi/4}$ & 010 \\
$180°$ & $e^{i\pi}$ & 110 \\
$225°$ & $e^{i5\pi/4}$ & 111 \\
$270°$ & $e^{i3\pi/2}$ & 101 \\
$315°$ & $e^{i7\pi/4}$ & 100 \\
\hline
\end{tabular}
\end{center}

Nota: Cada par de símbolos adyacentes (consecutivos en fase) tiene exactamente 1 bit diferente.
\end{tcolorbox}

\subsection{Densidad espectral de potencia}

$$S_x(\nu) = \frac{A^2\sigma_a^2}{4T_s}\left[|H(\nu-\nu_0)|^2 + |H(\nu+\nu_0)|^2\right]$$

El espectro está centrado en $\pm\nu_0$ y su forma depende de $|H(\nu)|^2$.

\section{Criterio de Nyquist y TSWOC - Evitar la interferencia}

\subsection{El problema: Interferencia entre símbolos (ISI)}

Cuando transmitimos símbolos consecutivos rápidamente, cada pulso $h(t)$ tiene una duración que se extiende más allá de su período $T_s$. Esto causa que las "colas" de un símbolo interfieran con los siguientes.

\textbf{Consecuencia:} Incluso sin ruido, la ISI degrada la detección y causa errores.

\textbf{Ejemplo visual:} Si $h(t)$ dura mucho tiempo, cuando muestreamos en $t = kT_s$ para detectar el símbolo $a_k$, también recibimos contribuciones de $a_{k-1}$, $a_{k+1}$, etc.

\textbf{Solución:} Diseñar $h(t)$ para que satisfaga el criterio TSWOC o Nyquist.

\subsection{Criterio TSWOC (Time-Shifted Waveform Orthogonality)}

\begin{tcolorbox}[colback=azulclaro!10,colframe=azuloscuro,title=Condición matemática]
Un filtro $h(t)$ satisface TSWOC si sus versiones desplazadas temporalmente son ortogonales:

$$\langle h(t), h(t-lT_s)\rangle = \int_{\mathbb{R}} h(t)h^*(t-lT_s)dt = \delta_{l,0} = \begin{cases}
1 & \text{si } l=0\\
0 & \text{si } l\neq 0
\end{cases}$$

\textbf{Interpretación física:}
\begin{itemize}
\item Cuando $l=0$: el pulso tiene producto interno 1 consigo mismo (normalización)
\item Cuando $l\neq 0$: no hay correlación entre el pulso y sus versiones desplazadas
\end{itemize}
\end{tcolorbox}

\textbf{Consecuencias del TSWOC:}
\begin{enumerate}
\item \textbf{Ausencia de ISI:} En los instantes de muestreo $t_0 + kT_s$, solo el símbolo $a_k$ contribuye a la muestra. Las contribuciones de todos los demás símbolos son cero.

\item \textbf{Detección símbolo por símbolo:} Cada símbolo puede decidirse independientemente sin necesidad de considerar los símbolos vecinos.

\item \textbf{Máximo SNR:} Cuando se usa con filtro adaptado en recepción, se maximiza la relación señal/ruido en el muestreador.
\end{enumerate}

\subsection{Criterio de Nyquist (equivalente en frecuencia)}

El criterio TSWOC en tiempo es equivalente al criterio de Nyquist en frecuencia:
$$\frac{1}{T_s}\sum_{p=-\infty}^{\infty} |H(\nu - \frac{p}{T_s})|^2 = 1, \quad \forall \nu \in \mathbb{R}$$

O en tiempo: $h(lT_s) = \delta_{l,0}$ (el pulso vale 1 en $t=0$ y 0 en todos los múltiplos de $T_s$)

\textbf{Interpretación:} Si sumamos todas las réplicas espectrales de $|H(\nu)|^2$ desplazadas por múltiplos de $1/T_s$, obtenemos una constante.

\textbf{Ancho de banda mínimo (Teorema):}
$$B_{min} = \frac{1}{2T_s} = \frac{D_s}{2}$$

Este es el ancho de banda teórico mínimo para transmitir a $D_s$ símbolos por segundo sin ISI. Corresponde al filtro rectangular ideal (no realizable en práctica porque no es causal y tiene duración infinita).

\subsection{Filtro Root Raised Cosine (RRC) - Solución práctica}

El filtro rectangular ideal no es realizable (infinita duración, no causal). El filtro más usado en la práctica es el \textbf{Root Raised Cosine (RRC)}, parametrizado por el factor de roll-off $\beta \in [0,1]$.

\textbf{Dominio temporal:}
$$h_{\beta}(t) = \frac{\sin[\pi(1-\beta)t/T_s] + \frac{4\beta t}{T_s}\cos[\pi(1+\beta)t/T_s]}{\pi t/T_s (1-16\beta^2t^2/T_s^2)}$$

\textbf{Dominio frecuencial:}
$$H_{\beta}(\nu) = \begin{cases}
\sqrt{T_s} & |\nu| \leq \frac{1-\beta}{2T_s}\\
\sqrt{\frac{T_s}{2}\left[1 + \cos\left(\frac{\pi T_s}{\beta}(|\nu| - \frac{1-\beta}{2T_s})\right)\right]} & \frac{1-\beta}{2T_s} < |\nu| \leq \frac{1+\beta}{2T_s}\\
0 & |\nu| > \frac{1+\beta}{2T_s}
\end{cases}$$

\textbf{Propiedades fundamentales:}
\begin{itemize}
\item \textbf{Propiedad clave:} $h_{\beta}(t) * h_{\beta}(-t)$ satisface el criterio de Nyquist

Esto significa: RRC en emisión $\times$ RRC en recepción = Raised Cosine (satisface Nyquist)

\item TSWOC: $\langle h_{\beta}(t), h_{\beta}(t-lT_s)\rangle = \delta_{l,0}$

\item Ancho de banda: $B = \frac{1+\beta}{2T_s}$

\item Exceso de banda: $\beta \times 100\%$ sobre el mínimo teórico $B_{min} = \frac{1}{2T_s}$

\item \textbf{Reparto emisor/receptor:} 
  \begin{itemize}
  \item Emisor: filtra con RRC
  \item Receptor: filtra con RRC (filtro adaptado)
  \item Cascada: $H_{RRC} \times H_{RRC} = H_{RC}$ (satisface Nyquist → sin ISI)
  \end{itemize}
\end{itemize}

\textbf{Parámetro $\beta$ (roll-off factor):}
\begin{itemize}
\item $\beta = 0$: Filtro rectangular (mínimo ancho de banda, NO realizable, oscilaciones infinitas)
\item $\beta = 1$: Filtro coseno completo (ancho de banda doble, fácil implementación, decaimiento rápido)
\item \textbf{Compromiso típico: $\beta = 0.25$ a $0.5$}
  \begin{itemize}
  \item Ancho de banda razonable
  \item Decaimiento temporal aceptable (menos oscilaciones)
  \item Sensibilidad moderada a errores de sincronización
  \end{itemize}
\end{itemize}

\textbf{¿Por qué "Root" (raíz)?}

Porque queremos que el producto de dos RRC dé un Raised Cosine:
$$H_{RRC}(\nu) \cdot H_{RRC}(\nu) = H_{RC}(\nu)$$
$$H_{RRC}(\nu) = \sqrt{H_{RC}(\nu)}$$

\textbf{Ventajas del esquema RRC-RRC:}
\begin{enumerate}
\item \textbf{Filtro adaptado óptimo:} RRC en recepción maximiza SNR
\item \textbf{Sin ISI:} La cascada satisface Nyquist
\item \textbf{Reparto simétrico:} Emisor y receptor comparten la tarea de filtrado
\end{enumerate}

\section{Canal AWGN y Receptor Óptimo}

\subsection{Modelo de canal AWGN - El peor enemigo}

El canal \textbf{AWGN (Additive White Gaussian Noise)} es el modelo más fundamental en comunicaciones:

$$r(t) = s(t) + n(t)$$

donde:
\begin{itemize}
\item $s(t)$: señal transmitida (conocida en emisión)
\item $n(t)$: ruido aditivo blanco gaussiano (aleatorio, impredecible)
\item $r(t)$: señal recibida (lo único que observamos)
\end{itemize}

\textbf{Características del ruido AWGN:}
\begin{itemize}
\item \textbf{Media:} $E[n(t)] = 0$ (no introduce sesgo sistemático)
\item \textbf{PSD (densidad espectral):} $S_n(\nu) = N_0$ (constante en todas las frecuencias → "blanco")
\item \textbf{Distribución:} $n(t) \sim \mathcal{N}(0, N_0)$ (gaussiana → campana de Gauss)
\item \textbf{Autocorrelación:} $R_n(\tau) = N_0\delta(\tau)$ (muestras independientes en instantes diferentes)
\end{itemize}

\textbf{¿Por qué es importante el modelo AWGN?}
\begin{enumerate}
\item Modelo fundamental del ruido térmico en electrónica
\item Ruido de fondo en comunicaciones inalámbricas
\item Aproximación válida en muchos canales reales
\item Permite análisis matemático riguroso
\item Si funciona bien con AWGN, funcionará razonablemente en otros canales
\end{enumerate}

\textbf{Parámetro $N_0$:} Potencia del ruido por Hz de ancho de banda. Mayor $N_0$ → peor canal.

\subsection{Criterio Maximum Likelihood (ML) - Decisión óptima}

El receptor debe decidir qué símbolo fue transmitido basándose únicamente en la señal recibida ruidosa $r(t)$.

\textbf{Criterio ML (máxima verosimilitud):}
$$\hat{k} = \arg\max_{k \in \{1,\ldots,M\}} P(r(t)|s^{(k)}(t) \text{ transmitido})$$

\textbf{Interpretación:} Elegir el símbolo que hace más probable la observación $r(t)$.

\textbf{Pregunta:} Dado que recibí $r(t)$, ¿cuál símbolo $s^{(k)}(t)$ es más probable que se haya transmitido?

\textbf{Para símbolos equiprobables} (cada símbolo tiene misma probabilidad $1/M$), el criterio ML se simplifica a:
$$\hat{k} = \arg\min_{k} ||r(t) - s^{(k)}(t)||^2 = \arg\min_{k} \int_{\mathbb{R}} |r(t) - s^{(k)}(t)|^2 dt$$

\textbf{Regla simple:} Elegir el símbolo más cercano (en distancia euclídea) a la señal recibida.

\textbf{Interpretación geométrica:}
\begin{itemize}
\item Las señales $s^{(k)}(t)$ son vectores en un espacio infinito-dimensional
\item La señal recibida $r(t) = s^{(k_0)}(t) + n(t)$ es el símbolo transmitido más ruido
\item El ruido "desplaza" $r(t)$ aleatoriamente del símbolo transmitido
\item Elegimos el símbolo más cercano a $r(t)$
\end{itemize}

\textbf{¿Por qué es óptimo?} El criterio ML minimiza la probabilidad de error cuando los símbolos son equiprobables.

\subsection{Estructura del receptor óptimo - Implementación práctica}

\begin{tcolorbox}[colback=azulclaro!10,colframe=azuloscuro,title=Receptor ML: Tres bloques fundamentales]
\textbf{1. Filtro adaptado (Matched Filter):} $g(t) = h^*(t_0 - t)$

\textbf{¿Qué hace?}
\begin{itemize}
\item Es la versión conjugada e invertida temporalmente del pulso de forma $h(t)$
\item \textbf{Maximiza la relación señal/ruido (SNR)} en el instante de muestreo
\item Con TSWOC: además elimina completamente la ISI
\end{itemize}

\textbf{¿Por qué es óptimo?} El teorema del filtro adaptado demuestra que $g(t) = h^*(-t)$ es el filtro lineal que maximiza SNR.

\textbf{Intuición:} El filtro "busca" en la señal recibida la forma del pulso transmitido. Cuando encuentra máxima correlación, produce un pico.

\textbf{2. Muestreador:} Muestrea en $t_n = t_0 + nT_s$

\textbf{¿Qué hace?}
\begin{itemize}
\item Convierte la señal continua en muestras discretas
\item $t_0$: instante óptimo de muestreo (máxima apertura del "ojo" en diagrama de ojo)
\item Una muestra por símbolo
\end{itemize}

\textbf{Salida si TSWOC se satisface:}
$$z_k = a_k + \gamma_k$$
donde:
\begin{itemize}
\item $a_k$: símbolo transmitido (señal útil)
\item $\gamma_k \sim \mathcal{N}(0, N_0)$: ruido gaussiano (único término de degradación)
\end{itemize}

\textbf{¡Logro importante!} La ISI se elimina completamente. Solo queda el ruido.

\textbf{3. Decisor (Detector):} Compara con regiones de decisión

\textbf{¿Qué hace?}
\begin{itemize}
\item Recibe la muestra compleja $z_k$
\item Determina a qué región de decisión pertenece
\item Decide qué símbolo $\hat{a}_k$ se transmitió
\end{itemize}

\textbf{Regla de decisión ML:}
$$\hat{d}_k = \arg\min_{d \in \mathcal{A}} |z_k - d|^2$$

Elegir el símbolo de la constelación $\mathcal{A}$ más cercano a $z_k$.

\textbf{Fronteras de decisión:} Son las bisectrices perpendiculares de los segmentos que unen símbolos vecinos (diagrama de Voronoi).

\textbf{Para modulación con portadora:}
\begin{itemize}
\item \textbf{Demodulación coherente} (requiere conocer fase $\theta_0$ de la portadora)
\item Dos ramas paralelas:
  \begin{itemize}
  \item Rama I (en fase): multiplica por $\cos(2\pi\nu_0 t + \theta_0)$
  \item Rama Q (cuadratura): multiplica por $-\sin(2\pi\nu_0 t + \theta_0)$
  \end{itemize}
\item Filtro adaptado en cada rama
\item Muestreo en cada rama → obtenemos $a_k$ y $b_k$
\item Decisión en plano complejo: $d_k = a_k + ib_k$
\end{itemize}

\textbf{Desafío práctico:} Recuperar la fase $\theta_0$ de la portadora (sincronización de portadora) y el instante óptimo $t_0$ (sincronización de símbolo).
\end{tcolorbox}

\subsection{Diagrama de ojo - Herramienta de diagnóstico visual}

El diagrama de ojo se obtiene superponiendo muchos períodos de símbolo de la señal recibida (antes del muestreador).

\textbf{¿Cómo se interpreta?}
\begin{itemize}
\item \textbf{Apertura vertical (altura del "ojo"):} Margen ante ruido
  \begin{itemize}
  \item Ojo bien abierto → mucho margen → baja tasa de error
  \item Ojo cerrado → poco margen → alta tasa de error
  \end{itemize}

\item \textbf{Apertura horizontal (ancho del "ojo"):} Sensibilidad al jitter temporal
  \begin{itemize}
  \item Ojo ancho → tolerante a errores de sincronización
  \item Ojo estrecho → muy sensible al instante de muestreo
  \end{itemize}

\item \textbf{Centro del ojo:} Instante óptimo de muestreo $t_0$
  \begin{itemize}
  \item Punto de máxima apertura
  \item Mínima ISI y máximo SNR
  \end{itemize}

\item \textbf{Líneas horizontales:} Niveles de decisión (umbrales)
  \begin{itemize}
  \item Separan las regiones de decisión
  \item Idealmente en la mitad de la apertura vertical
  \end{itemize}

\item \textbf{Grosor de las trazas:} Cantidad de ruido e ISI
  \begin{itemize}
  \item Trazas finas → poco ruido/ISI
  \item Trazas gruesas → mucho ruido/ISI
  \end{itemize}
\end{itemize}

\textbf{Efecto del roll-off $\beta$:}
\begin{itemize}
\item $\beta$ pequeño (→ 0): Mayor eficiencia espectral, pero mayores oscilaciones temporales → ojo menos abierto, más sensible a timing
\item $\beta$ grande (→ 1): Menor eficiencia espectral, pero pulsos más "suaves" → ojo más abierto, menos sensible a timing
\end{itemize}

\section{Análisis de Prestaciones - ¿Qué tan bien funciona?}

\subsection{SNR - Medida de calidad fundamental}

La relación señal/ruido (Signal-to-Noise Ratio) mide la calidad del canal.

\textbf{SNR por símbolo:} 
$$\text{SNR}_s = \frac{E_s}{N_0}$$

donde $E_s$ es la energía promedio por símbolo y $N_0$ es la densidad espectral del ruido.

\textbf{SNR por bit:} 
$$\frac{E_b}{N_0} = \frac{E_s}{\log_2(M) \cdot N_0} = \frac{\text{SNR}_s}{\log_2(M)}$$

\textbf{¿Por qué $E_b/N_0$ es fundamental?}
\begin{itemize}
\item Es independiente del ancho de banda usado
\item Permite comparar diferentes esquemas de modulación
\item Mide eficiencia energética pura
\item Umbral para comunicación confiable (Límite de Shannon)
\end{itemize}

\textbf{Conversión a dB:} $(E_b/N_0)_{dB} = 10\log_{10}(E_b/N_0)$

\textbf{Ejemplo:} $E_b/N_0 = 10 \Rightarrow (E_b/N_0)_{dB} = 10$ dB

\textbf{Valores típicos necesarios:}
\begin{itemize}
\item BPSK (BER=$10^{-5}$): $E_b/N_0 \approx 9.6$ dB
\item QPSK (BER=$10^{-5}$): $E_b/N_0 \approx 9.6$ dB (igual que BPSK)
\item 16-QAM (BER=$10^{-5}$): $E_b/N_0 \approx 14$ dB
\item 64-QAM (BER=$10^{-5}$): $E_b/N_0 \approx 18$ dB
\end{itemize}

\textbf{Interpretación:} Mayor $M$ requiere mayor SNR para misma tasa de error.

\subsection{Probabilidad de error - Cuantificar errores}

\textbf{Función erfc (complementary error function):}
$$\text{erfc}(x) = \frac{2}{\sqrt{\pi}}\int_{x}^{\infty}e^{-t^2}dt$$

Está relacionada con la distribución gaussiana. Decrece muy rápidamente con $x$.

\textbf{Función Q (cola de la gaussiana):}
$$Q(x) = \frac{1}{2}\text{erfc}\left(\frac{x}{\sqrt{2}}\right) = P(X > x) \text{ para } X \sim \mathcal{N}(0,1)$$

$Q(x)$ es la probabilidad de que una variable aleatoria gaussiana estándar exceda el valor $x$.

\textbf{Union Bound (cota superior de la probabilidad de error):}
$$P_{se} \leq \frac{1}{M}\sum_{k=1}^{M}\sum_{l\neq k} \frac{1}{2}\text{erfc}\left(\frac{||s^{(k)} - s^{(l)}||}{2\sqrt{2N_0}}\right)$$

\textbf{¿Qué dice?} La probabilidad de error de símbolo es menor o igual que la suma de las probabilidades de confundir cada par de símbolos.

\textbf{Aproximación para alto SNR (la más útil):}
$$P_{se} \approx \frac{N_{min}}{2}\text{erfc}\left(\sqrt{\frac{d_{min}^2}{8N_0}}\right)$$

donde:
\begin{itemize}
\item $d_{min}$: distancia mínima entre símbolos de la constelación
\item $N_{min}$: número máximo de vecinos a distancia $d_{min}$
\end{itemize}

\textbf{Interpretación clave:}
\begin{itemize}
\item La tasa de error depende exponencialmente de $d_{min}^2/(N_0)$
\item Doblar $d_{min}$ reduce $P_{se}$ en varios órdenes de magnitud
\item Por eso el diseño de constelaciones busca maximizar $d_{min}$
\end{itemize}

\subsection{Modulaciones específicas (TSWOC)}

\textbf{M-PAM:}
$$P_{se} \approx \frac{M-1}{M}\text{erfc}\left(\sqrt{\frac{3E_{avg}}{(M^2-1)N_0}}\right)$$

\textbf{M-PSK:}
$$P_{se} \approx \text{erfc}\left(\sqrt{\frac{E_s}{N_0}}\sin^2\left(\frac{\pi}{M}\right)\right)$$

\textbf{M-QAM:}
$$P_{se} \approx 2\left(1-\frac{1}{\sqrt{M}}\right)\text{erfc}\left(\sqrt{\frac{3E_{avg}}{2(M-1)N_0}}\right)$$

\textbf{Con Gray:} $P_{be} \approx \frac{P_{se}}{\log_2(M)}$

\subsection{Eficiencia de potencia}

$$\rho = \frac{d_{min}^2}{E_{avg}}$$

Mide eficiencia en uso de energía. Mayor $\rho$ → mejor.

Para $M$ grande: QAM > PSK > PAM

\subsection{Diseño óptimo de constelaciones}

\begin{tcolorbox}[colback=azulclaro!10,colframe=azuloscuro,title=Principios]
Para minimizar $P_e$ con $E_{avg}$ fija:
\begin{enumerate}
\item \textbf{Maximizar $d_{min}$} (impacto dominante, $P_e \propto e^{-d_{min}^2}$)
\item \textbf{Minimizar número de vecinos} (impacto secundario)
\end{enumerate}

\textbf{Comparación:}
\begin{itemize}
\item QAM: mejor para $M$ grande (mayor $d_{min}$ con misma $E_{avg}$)
\item PSK: PAPR=1, mejor para amplificadores no lineales
\end{itemize}
\end{tcolorbox}

\section{Codificación de Canal - Añadir redundancia inteligente}

\subsection{Motivación profunda - ¿Por qué necesitamos códigos?}

\begin{tcolorbox}[colback=rojoclaro!10,colframe=rojoclaro,title=🧠 RECORDATORIO FUNDAMENTAL]
\textbf{¿Por qué falla la comunicación sin códigos?}

Cuando transmitimos bits por un canal real:
\begin{itemize}
\item El \textbf{ruido AWGN} hace que la señal recibida $r = s + n$ sea aleatoria
\item Incluso con SNR alto, siempre hay una probabilidad $P_e > 0$ de decidir mal
\item Sin protección, un solo bit erróneo puede corromper una imagen, archivo, mensaje
\end{itemize}

\textbf{Esto era así porque:} La decisión símbolo por símbolo trata cada bit de forma \emph{independiente}. No aprovechamos que conocemos la \emph{estructura estadística} del mensaje completo.
\end{tcolorbox}

\textbf{Problema fundamental:} El canal introduce errores aleatorios. Para un canal BSC (Binary Symmetric Channel) con probabilidad de error $p$:
\begin{itemize}
\item Bit transmitido 0 → recibido 0 con probabilidad $(1-p)$, recibido 1 con probabilidad $p$
\item Bit transmitido 1 → recibido 1 con probabilidad $(1-p)$, recibido 0 con probabilidad $p$
\end{itemize}

Sin códigos, si $p = 10^{-3}$, en 1000 bits tendremos aproximadamente 1 error. ¡Inaceptable para muchas aplicaciones!

\textbf{Idea brillante de la codificación:} Añadir redundancia \emph{controlada} (bits extra calculados matemáticamente) para:
\begin{enumerate}
\item \textbf{Detectar} errores: saber que algo salió mal
\item \textbf{Localizar} errores: saber \emph{dónde} ocurrieron
\item \textbf{Corregir} errores: recuperar el mensaje original sin retransmisión
\end{enumerate}

\textbf{Compromiso fundamental (trade-off inevitable):}
\begin{itemize}
\item \textbf{Ventaja:} Reduce drásticamente la tasa de error final
  \begin{itemize}
  \item Puede bajar de $P_e = 10^{-3}$ a $P_e = 10^{-9}$ o menos
  \item Permite aumentar alcance (menor potencia necesaria)
  \item Mejora robustez en canales difíciles (móvil, espacio profundo)
  \end{itemize}
\item \textbf{Desventaja:} Reduce el débito útil efectivo
  \begin{itemize}
  \item Transmitimos $n$ bits pero solo $k < n$ son información
  \item Tasa del código: $R_c = k/n < 1$
  \item Overhead: $(1 - R_c) \times 100\%$ (por ejemplo, 50\% para $R_c=1/2$)
  \end{itemize}
\end{itemize}

\textbf{Ejemplo pedagógico - Código de repetición (3,1):}
\begin{itemize}
\item Entrada: 1 bit de información (0 o 1)
\item Codificación: Repetir 3 veces
  \begin{itemize}
  \item Bit 0 → transmitir \textbf{000}
  \item Bit 1 → transmitir \textbf{111}
  \end{itemize}
\item Decodificación: Regla de mayoría
  \begin{itemize}
  \item Recibimos 010 → mayoría es 0 → decidimos 0
  \item Recibimos 101 → mayoría es 1 → decidimos 1
  \item Recibimos 011 → mayoría es 1 → decidimos 1
  \end{itemize}
\item Capacidad: Corrige hasta 1 error de los 3 bits
\item Tasa: $R_c = 1/3 \approx 0.33$ (solo 33\% de eficiencia, muy ineficiente pero muy robusto)
\end{itemize}

\begin{tcolorbox}[colback=verdeclaro!10,colframe=verdeclaro,title=⚡ TEOREMA DE SHANNON - Resultado más importante]
\textbf{Segundo Teorema de Shannon (Teorema de Codificación de Canal):}

Existe un límite fundamental llamado \textbf{capacidad del canal} $C$ tal que:
\begin{itemize}
\item Si $R < C$ (tasa de información menor que capacidad): \textbf{Existen códigos} que permiten comunicar con probabilidad de error \textbf{arbitrariamente pequeña} ($P_e \to 0$)
\item Si $R > C$: \textbf{Imposible} comunicar de forma confiable, sin importar qué código usemos
\end{itemize}

\textbf{Para canal AWGN con ancho de banda $B$ y SNR = $S/N$:}
$$\boxed{C = B\log_2\left(1 + \frac{S}{N}\right) \quad \text{[bits/s]}}$$

\textbf{Forma alternativa (normalizada por ancho de banda):}
$$\frac{C}{B} = \log_2\left(1 + \frac{E_b R}{N_0 B}\right) = \log_2\left(1 + \frac{E_b}{N_0} \cdot \frac{R}{B}\right)$$

\textbf{Límite de Shannon ($R/B \to 0$, ancho de banda infinito):}
$$\frac{E_b}{N_0}\bigg|_{\text{Shannon}} = \ln(2) \approx 0.693 \approx -1.59 \text{ dB}$$

Este es el \emph{límite teórico absoluto}. Ningún sistema real puede operar por debajo de $E_b/N_0 = -1.59$ dB.

\textbf{Interpretación práctica:}
\begin{itemize}
\item Sistemas sin código: necesitan típicamente $E_b/N_0 \approx 10$ dB para BER = $10^{-5}$
\item Códigos modernos (Turbo, LDPC): operan a $\approx 1$-2 dB del límite de Shannon
\item \textbf{Ganancia de codificación:} $\approx 8$-9 dB (¡factor 6-8 en potencia!)
\end{itemize}
\end{tcolorbox}

\textbf{¿Qué nos dice el teorema de Shannon?}
\begin{enumerate}
\item La comunicación confiable es \textbf{posible} incluso con ruido (no necesitamos SNR infinito)
\item Hay un límite fundamental que no podemos superar (la física lo impide)
\item Existe un trade-off entre tasa de información y robustez
\item El \textbf{ancho de banda} y la \textbf{potencia} son intercambiables (hasta cierto punto)
\end{enumerate}

\textbf{Tipos de estrategias de codificación:}
\begin{itemize}
\item \textbf{ARQ (Automatic Repeat reQuest):} Solo detectar errores y pedir retransmisión
  \begin{itemize}
  \item Ventaja: Simple, usa menos redundancia
  \item Desventaja: Requiere canal de retorno, aumenta latencia
  \item Uso: TCP/IP, protocolos de red
  \end{itemize}
\item \textbf{FEC (Forward Error Correction):} Corregir errores automáticamente
  \begin{itemize}
  \item Ventaja: No necesita retransmisión, latencia fija
  \item Desventaja: Más redundancia necesaria
  \item Uso: Broadcasting (TV, radio), tiempo real (voz), comunicaciones espaciales
  \end{itemize}
\item \textbf{Híbrido (HARQ):} Combinar ambas estrategias
  \begin{itemize}
  \item FEC para errores comunes, ARQ para errores raros
  \item Uso: LTE, 5G, WiFi moderno
  \end{itemize}
\end{itemize}

\subsection{Códigos de bloque lineales C(n,k) - Estructura algebraica}

\begin{tcolorbox}[colback=azulclaro!10,colframe=azuloscuro,title=🧠 FUNDAMENTO: ¿Por qué "lineales"?]
\textbf{Esto era así porque:} Trabajamos en el cuerpo finito $\mathbb{F}_2 = \{0,1\}$ (aritmética módulo 2).

\textbf{Operaciones en $\mathbb{F}_2$:}
\begin{itemize}
\item Suma: $0 + 0 = 0$, $0 + 1 = 1$, $1 + 0 = 1$, $1 + 1 = 0$ (XOR)
\item Multiplicación: $0 \cdot 0 = 0$, $0 \cdot 1 = 0$, $1 \cdot 0 = 0$, $1 \cdot 1 = 1$ (AND)
\end{itemize}

\textbf{Propiedad "lineal":} Si $\mathbf{c}_1$ y $\mathbf{c}_2$ son palabras código válidas, entonces $\mathbf{c}_1 + \mathbf{c}_2$ (suma módulo 2, bit a bit) también es palabra código válida.

\textbf{Consecuencias poderosas:}
\begin{enumerate}
\item La palabra todo-ceros $\mathbf{0}$ siempre pertenece al código
\item El código forma un \textbf{subespacio vectorial} de $\mathbb{F}_2^n$
\item Podemos usar \textbf{álgebra lineal} para codificar y decodificar
\item La estructura algebraica simplifica el análisis y la implementación
\end{enumerate}
\end{tcolorbox}

Un código de bloque lineal $\mathcal{C}(n,k)$ transforma bloques de $k$ bits de información en bloques de $n$ bits codificados, añadiendo $n-k$ bits de paridad.

\textbf{Parámetros fundamentales:}
\begin{itemize}
\item $n$: longitud de la palabra código (bits transmitidos en total)
\item $k$: longitud del mensaje de información (bits útiles de datos)
\item $M = 2^k$: número de palabras código válidas (de $2^n$ posibles en total)
  \begin{itemize}
  \item Solo $2^k$ de las $2^n$ secuencias posibles son palabras código válidas
  \item Esto es clave: la \textbf{redundancia} viene de usar solo un subconjunto
  \end{itemize}
\item \textbf{Tasa del código (code rate):} $R_c = k/n$ (fracción de bits útiles)
  \begin{itemize}
  \item $R_c = 1$: Sin código (todos los bits son información)
  \item $R_c \to 0$: Mucha redundancia (muy protegido pero ineficiente)
  \item Típicamente: $R_c \in [1/3, 9/10]$
  \end{itemize}
\item Bits de paridad: $n - k$ (redundancia añadida)
\item \textbf{Overhead:} $(1 - R_c) = (n-k)/n$ (fracción de bits "desperdiciados")
\end{itemize}

\textbf{Ejemplo detallado - Código Hamming (7,4):}
\begin{itemize}
\item 4 bits de información: $\mathbf{m} = [m_0, m_1, m_2, m_3]$
\item 3 bits de paridad: calculados a partir de $\mathbf{m}$
\item 7 bits totales transmitidos: $\mathbf{c} = [m_0, m_1, m_2, m_3, p_0, p_1, p_2]$
\item $R_c = 4/7 \approx 0.571$ (57.1\% eficiencia)
\item Overhead = 3/7 ≈ 43\% (casi la mitad son bits de paridad)
\item Puede corregir \textbf{1 error} en cualquier posición
\item Si transmitimos 1 Mbps codificado → débito útil = 571 kbps
\end{itemize}

\textbf{Matriz generadora $\mathbf{G}$ (forma sistemática):}
$$\mathbf{G} = [\mathbf{I}_k | \mathbf{P}]_{k \times n}$$

donde:
\begin{itemize}
\item $\mathbf{I}_k$: matriz identidad $k \times k$ (preserva bits de información)
\item $\mathbf{P}$: matriz de paridad $k \times (n-k)$ (genera redundancia)
\item Dimensión total: $k$ filas, $n$ columnas
\end{itemize}

\textbf{Codificación (multiplicación matricial módulo 2):}
$$\boxed{\mathbf{c} = \mathbf{m}\mathbf{G} \quad (\text{mod } 2)}$$

donde:
\begin{itemize}
\item $\mathbf{m} = [m_0, m_1, \ldots, m_{k-1}]$: mensaje de información (vector fila $1 \times k$)
\item $\mathbf{c} = [c_0, c_1, \ldots, c_{n-1}]$: palabra código (vector fila $1 \times n$)
\item Operación: producto matriz-vector con aritmética módulo 2
\end{itemize}

\textbf{Forma sistemática (ventaja práctica):}
$$\mathbf{c} = [\underbrace{m_0, m_1, \ldots, m_{k-1}}_{\text{Información (visible)}}, \underbrace{p_0, p_1, \ldots, p_{n-k-1}}_{\text{Paridad (calculada)}}]$$

Los primeros $k$ bits de $\mathbf{c}$ son \textbf{exactamente} $\mathbf{m}$ (información transparente). Esto simplifica la implementación:
\begin{itemize}
\item El codificador solo necesita calcular los bits de paridad
\item El decodificador puede leer directamente la información (si no hay errores)
\item Facilita interoperabilidad con sistemas sin código
\end{itemize}

\textbf{Ejemplo concreto - Hamming (7,4):}
$$\mathbf{G} = \begin{pmatrix}
1 & 0 & 0 & 0 & 1 & 1 & 0\\
0 & 1 & 0 & 0 & 1 & 0 & 1\\
0 & 0 & 1 & 0 & 0 & 1 & 1\\
0 & 0 & 0 & 1 & 1 & 1 & 1
\end{pmatrix}$$

Si $\mathbf{m} = [1, 0, 1, 1]$:
$$\mathbf{c} = [1, 0, 1, 1] \cdot \mathbf{G} = [1, 0, 1, 1, 0, 0, 0]$$

\textbf{Matriz de paridad $\mathbf{H}$ (verificación):}
$$\mathbf{H} = [\mathbf{P}^T | \mathbf{I}_{n-k}]_{(n-k) \times n}$$

Dimensión: $(n-k)$ filas, $n$ columnas

\textbf{Propiedad fundamental (ortogonalidad):}
$$\boxed{\mathbf{c}\mathbf{H}^T = \mathbf{0} \quad \text{si y solo si } \mathbf{c} \text{ es palabra código válida}}$$

\textbf{¿Por qué funciona esto?} 
\begin{itemize}
\item $\mathbf{G}$ genera el subespacio código (filas son base)
\item $\mathbf{H}$ define el complemento ortogonal (verifica pertenencia)
\item Se cumple siempre: $\mathbf{G}\mathbf{H}^T = \mathbf{0}$ (ortogonalidad algebraica)
\item Por tanto: $\mathbf{c}\mathbf{H}^T = (\mathbf{m}\mathbf{G})\mathbf{H}^T = \mathbf{m}(\mathbf{G}\mathbf{H}^T) = \mathbf{0}$
\end{itemize}

\textbf{¿Para qué sirve $\mathbf{H}$? (usos prácticos):}
\begin{enumerate}
\item \textbf{Verificar validez:} Calcular $\mathbf{r}\mathbf{H}^T$ donde $\mathbf{r}$ es palabra recibida
  \begin{itemize}
  \item Si $\mathbf{r}\mathbf{H}^T = \mathbf{0}$ → sin errores detectados (o error no detectable)
  \item Si $\mathbf{r}\mathbf{H}^T \neq \mathbf{0}$ → error detectado con certeza
  \end{itemize}
\item \textbf{Localizar errores:} El síndrome $\mathbf{s} = \mathbf{r}\mathbf{H}^T$ indica \emph{dónde} está el error
\item \textbf{Corregir errores:} Usar tabla de síndromes o algoritmos para encontrar el patrón de error
\end{enumerate}

\textbf{Complejidad computacional:}
\begin{itemize}
\item Codificación: $O(kn)$ operaciones XOR (rápido, $\approx$ microsegundos)
\item Decodificación: Depende del algoritmo
  \begin{itemize}
  \item Tabla completa: $O(1)$ tiempo, $O(2^{n-k})$ memoria (imposible para $n$ grande)
  \item Algoritmos iterativos (LDPC, Turbo): $O(n \cdot I_{max})$ donde $I_{max}$ es número de iteraciones
  \end{itemize}
\end{itemize}

\subsection{Distancia de Hamming - Métrica fundamental del espacio binario}

\begin{tcolorbox}[colback=rojoclaro!10,colframe=rojoclaro,title=🧠 FUNDAMENTO GEOMÉTRICO]
\textbf{Esto era así porque:} Las palabras binarias de longitud $n$ forman un espacio discreto de $2^n$ puntos (el hipercubo binario $\mathbb{F}_2^n$).

\textbf{Geometría del espacio binario:}
\begin{itemize}
\item Cada palabra es un vértice del hipercubo $n$-dimensional
\item La distancia de Hamming mide el número de aristas entre dos vértices
\item Equivalentemente: número de coordenadas que difieren
\end{itemize}

\textbf{Ejemplo visual ($n=3$, cubo 3D):}
\begin{verbatim}
        011 ----------- 111
       /|              /|
      / |             / |
    001----------101  |
     |  |          |  |
     | 010---------|--110
     | /           | /
     |/            |/
    000-----------100
\end{verbatim}

\textbf{Distancias desde 000:}
\begin{itemize}
\item $d_H(000, 001) = 1$ (un bit diferente, vecinos adyacentes)
\item $d_H(000, 011) = 2$ (dos bits diferentes, distancia 2)
\item $d_H(000, 111) = 3$ (tres bits diferentes, vértices opuestos)
\end{itemize}

\textbf{Propiedad clave:} En cada vértice hay exactamente $n$ vecinos a distancia 1 (uno por cada bit que podemos invertir).
\end{tcolorbox}

La distancia de Hamming es la \textbf{métrica natural} para medir diferencias entre secuencias binarias.

\textbf{Peso de Hamming (Hamming weight):}
$$w_H(\mathbf{x}) = \text{número de bits iguales a 1 en } \mathbf{x}$$

\textbf{Ejemplos:}
\begin{itemize}
\item $w_H(0000000) = 0$ (palabra todo-ceros)
\item $w_H(1010110) = 4$ (cuatro unos)
\item $w_H(1111111) = 7$ (palabra todo-unos)
\item $w_H(1000000) = 1$ (peso mínimo no nulo)
\end{itemize}

\textbf{Distancia de Hamming entre dos palabras:}
$$\boxed{d_H(\mathbf{x},\mathbf{y}) = w_H(\mathbf{x} + \mathbf{y}) = w_H(\mathbf{x} \oplus \mathbf{y})}$$

donde $+$ es suma módulo 2 (XOR bit a bit).

\textbf{¿Por qué funciona esta definición?}
\begin{itemize}
\item $\mathbf{x} \oplus \mathbf{y}$ tiene un 1 exactamente en las posiciones donde $\mathbf{x}$ y $\mathbf{y}$ difieren
\item El peso de esta diferencia es el número de posiciones diferentes
\item Por tanto: $d_H(\mathbf{x},\mathbf{y})$ = número de bits diferentes
\end{itemize}

\textbf{Ejemplo detallado paso a paso:}
\begin{align*}
\mathbf{x} &= 1011010 \\
\mathbf{y} &= 1001110 \\
\hline
\mathbf{x} \oplus \mathbf{y} &= 0010100 \quad \text{(XOR bit a bit)}\\
d_H(\mathbf{x},\mathbf{y}) &= w_H(0010100) = 2
\end{align*}

\textbf{Verificación directa:} Comparando posición por posición:
\begin{itemize}
\item Posiciones 0,1,3,5,6: iguales ($x_i = y_i$)
\item Posiciones 2,4: diferentes ($x_2 \neq y_2$, $x_4 \neq y_4$)
\item Total: 2 diferencias → $d_H = 2$ ✓
\end{itemize}

\textbf{Propiedades matemáticas (es una métrica):}
\begin{enumerate}
\item \textbf{No negatividad:} $d_H(\mathbf{x},\mathbf{y}) \geq 0$
\item \textbf{Identidad:} $d_H(\mathbf{x},\mathbf{y}) = 0 \Leftrightarrow \mathbf{x} = \mathbf{y}$
\item \textbf{Simetría:} $d_H(\mathbf{x},\mathbf{y}) = d_H(\mathbf{y},\mathbf{x})$
\item \textbf{Desigualdad triangular:} $d_H(\mathbf{x},\mathbf{z}) \leq d_H(\mathbf{x},\mathbf{y}) + d_H(\mathbf{y},\mathbf{z})$
\end{enumerate}

\textbf{Para códigos lineales (propiedad especial):}
La distancia entre dos palabras código es igual al peso de su diferencia:
$$d_H(\mathbf{c}_i, \mathbf{c}_j) = w_H(\mathbf{c}_i + \mathbf{c}_j) = w_H(\mathbf{c}_k)$$

donde $\mathbf{c}_k = \mathbf{c}_i + \mathbf{c}_j$ es también palabra código (por linealidad).

\textbf{Consecuencia práctica importante:} Para calcular $d_{min}$ de un código lineal, basta encontrar la palabra código no nula de menor peso.

\textbf{Distancia mínima del código $d_{min}$ (parámetro crítico):}
$$\boxed{d_{min} = \min_{\mathbf{c}_i \neq \mathbf{c}_j} d_H(\mathbf{c}_i, \mathbf{c}_j)}$$

Es la \textbf{menor distancia} entre cualquier par de palabras código válidas diferentes.

\textbf{¿Por qué es fundamental $d_{min}$?}

Determina \textbf{completamente} la capacidad de detección y corrección del código:
\begin{itemize}
\item Mayor $d_{min}$ → palabras código más "separadas" → más robusto
\item Menor $d_{min}$ → palabras código más "cercanas" → más vulnerable
\end{itemize}

\textbf{Cálculo de $d_{min}$ (método directo):}
\begin{enumerate}
\item Generar todas las $2^k$ palabras código
\item Calcular distancias entre todos los pares: $\binom{2^k}{2}$ distancias
\item Tomar el mínimo
\end{enumerate}

\textbf{Problema:} Complejidad explosiva para $k$ grande ($O(2^{2k})$).

\textbf{Método eficiente (códigos lineales):}
\begin{enumerate}
\item Generar todas las $2^k$ palabras código
\item Calcular el peso de cada una: $w_H(\mathbf{c}_i)$
\item Tomar el peso mínimo no nulo: $d_{min} = \min_{\mathbf{c} \neq \mathbf{0}} w_H(\mathbf{c})$
\end{enumerate}

Complejidad: $O(2^k \cdot n)$ (mucho mejor).

\textbf{Método usando matriz $\mathbf{H}$ (más avanzado):}
$$d_{min} = \text{número mínimo de columnas linealmente dependientes de } \mathbf{H}$$

\textbf{Ejemplos con códigos conocidos:}
\begin{itemize}
\item Código de repetición $(n,1)$: $d_{min} = n$ (máxima separación posible)
  \begin{itemize}
  \item Palabras: 000...0 y 111...1
  \item Distancia: $n$ (todos los bits difieren)
  \end{itemize}
\item Código de paridad $(n, n-1)$: $d_{min} = 2$ (separación mínima)
  \begin{itemize}
  \item Cualquier cambio de 1 bit cambia paridad
  \item Pero cambio de 2 bits puede ser indetectable
  \end{itemize}
\item Hamming $(7,4)$: $d_{min} = 3$
  \begin{itemize}
  \item Palabras código separadas por al menos 3 bits
  \item Ejemplo: 0000000, 1101000, 0110100 tienen distancias ≥ 3
  \end{itemize}
\item Golay $(23,12)$: $d_{min} = 7$ (código muy potente)
\item Reed-Solomon: $d_{min} = n - k + 1$ (alcanza cota de Singleton)
\end{itemize}

\subsection{Capacidad de detección y corrección - Teorema fundamental}

\begin{tcolorbox}[colback=azulclaro!10,colframe=azuloscuro,title=⚡ TEOREMA FUNDAMENTAL DE LA CODIFICACIÓN]
Un código con distancia mínima $d_{min}$ tiene las siguientes capacidades garantizadas:

\textbf{Detectar:} hasta $\boxed{d_{min} - 1}$ errores

\textbf{Corregir:} hasta $\boxed{t = \left\lfloor \frac{d_{min}-1}{2} \right\rfloor}$ errores

\textbf{Detectar Y corregir simultáneamente:}
\begin{itemize}
\item Corregir $t$ errores Y detectar $e$ errores adicionales si $d_{min} \geq t + e + 1$
\end{itemize}
\end{tcolorbox}

\textbf{Demostración intuitiva (detección):}

\begin{tcolorbox}[colback=verdeclaro!10,colframe=verdeclaro,title=🔍 Demostración - Detección]
\textbf{Esto era así porque:} Las palabras código válidas están separadas por al menos $d_{min}$ bits.

\textbf{Razonamiento:}
\begin{enumerate}
\item Transmitimos una palabra código válida $\mathbf{c}_i$
\item Ocurren hasta $d_{min}-1$ errores → recibimos $\mathbf{r}$ con $d_H(\mathbf{c}_i, \mathbf{r}) \leq d_{min}-1$
\item \textbf{Pregunta clave:} ¿Puede $\mathbf{r}$ ser otra palabra código válida $\mathbf{c}_j$?
\item \textbf{Respuesta:} NO, porque $d_H(\mathbf{c}_i, \mathbf{c}_j) \geq d_{min}$ para todo $j \neq i$
\item Si solo cambiamos $d_{min}-1$ bits de $\mathbf{c}_i$, NO podemos llegar a ningún otro $\mathbf{c}_j$
\item Por tanto: $\mathbf{r}$ NO es palabra código válida → \textbf{detectamos error}
\end{enumerate}

\textbf{Ejemplo visual ($d_{min}=4$):}
\begin{verbatim}
 Círc. 000: palabra código
 . : palabra inválida

   Radio 1    Radio 2    Radio 3
   (1 error) (2 err)   (3 err)
      .         .         .
     . .       . .       . .
    . 000 .   . . .     . . .
     . .       . .       . .
      .         .         .

 Otra palabra código debe estar
 a distancia ≥ 4, fuera de estos círculos
\end{verbatim}

\textbf{Conclusión:} Con $d_{min}=4$, detectamos hasta 3 errores garantizado.
\end{tcolorbox}

\textbf{Demostración intuitiva (corrección):}

\begin{tcolorbox}[colback=azulclaro!10,colframe=azuloscuro,title=🔧 Demostración - Corrección]
\textbf{Idea geométrica:} Imaginemos "esferas de corrección" de radio $t$ alrededor de cada palabra código.

\textbf{Condición para corrección sin ambigüedad:}
Las esferas NO deben solaparse → necesitamos $d_{min} \geq 2t + 1$

\textbf{¿Por qué $2t+1$?}
\begin{itemize}
\item Esfera alrededor de $\mathbf{c}_i$ tiene radio $t$
\item Esfera alrededor de $\mathbf{c}_j$ tiene radio $t$
\item Para que no se toquen: distancia entre centros $\geq t + t + 1 = 2t + 1$
\item El +1 garantiza que no hay punto intermedio equidistante
\end{itemize}

\textbf{Visualización ($d_{min}=5$, $t=2$):}
\begin{verbatim}
   Esfera radio 2       Esfera radio 2
   alrededor C1         alrededor C2
       . .                 . .
      . . .               . . .
     .  C1  .           .  C2  .
      . . .               . . .
       . .                 . .
   ├─────┤           ├─────┤
      2 bits   1 bit      2 bits
               (gap)
   ├──────────────────────┤
        d_min = 5 bits
\end{verbatim}

\textbf{Razonamiento paso a paso:}
\begin{enumerate}
\item Transmitimos $\mathbf{c}_i$, ocurren $\leq t$ errores
\item Recibimos $\mathbf{r}$ con $d_H(\mathbf{c}_i, \mathbf{r}) \leq t$
\item $\mathbf{r}$ está dentro de la esfera de radio $t$ alrededor de $\mathbf{c}_i$
\item \textbf{Pregunta:} ¿Puede $\mathbf{r}$ estar dentro de otra esfera (de $\mathbf{c}_j$)?
\item Si $d_{min} \geq 2t+1$, las esferas no se solapan → \textbf{NO}
\item Por tanto: $\mathbf{c}_i$ es la única palabra código a distancia $\leq t$ de $\mathbf{r}$
\item \textbf{Decodificación:} Elegir $\mathbf{c}_i$ (palabra más cercana) → corrección exitosa
\end{enumerate}
\end{tcolorbox}

\textbf{Fórmula para capacidad de corrección:}
$$t = \left\lfloor \frac{d_{min}-1}{2} \right\rfloor$$

\textbf{¿Por qué el piso (floor)?}
\begin{itemize}
\item Si $d_{min}$ es impar: $d_{min} = 2t+1$ → $t = (d_{min}-1)/2$ (exacto)
\item Si $d_{min}$ es par: $d_{min} = 2t$ o $2t+1$ (no se usa completamente la capacidad)
\end{itemize}

\textbf{Ejemplos numéricos detallados:}

\textbf{1) Código con $d_{min} = 3$:}
\begin{itemize}
\item Detecta: hasta $3-1 = 2$ errores
\item Corrige: $t = \lfloor (3-1)/2 \rfloor = \lfloor 1 \rfloor = 1$ error
\item \textbf{Significado práctico:}
  \begin{itemize}
  \item Si 1 bit se corrompe → corregimos automáticamente
  \item Si 2 bits se corrompen → detectamos pero NO podemos corregir (ambiguo)
  \item Si 3 bits se corrompen → puede ser indetectable (llegamos a otra palabra código)
  \end{itemize}
\item \textbf{Ejemplo:} Código Hamming $(7,4)$
\end{itemize}

\textbf{2) Código con $d_{min} = 5$:}
\begin{itemize}
\item Detecta: hasta $5-1 = 4$ errores
\item Corrige: $t = \lfloor 4/2 \rfloor = 2$ errores
\item \textbf{Capacidades:}
  \begin{itemize}
  \item Hasta 2 errores → corrección automática
  \item 3-4 errores → detección sin corrección
  \item 5+ errores → puede fallar completamente
  \end{itemize}
\item \textbf{Ejemplo:} Código BCH $(31,16)$
\end{itemize}

\textbf{3) Código con $d_{min} = 7$:}
\begin{itemize}
\item Detecta: hasta 6 errores
\item Corrige: $t = \lfloor 6/2 \rfloor = 3$ errores
\item \textbf{Ejemplo:} Código Golay $(23,12)$
\end{itemize}

\textbf{4) Código de repetición $(n,1)$ con $d_{min} = n$:}
\begin{itemize}
\item Detecta: hasta $n-1$ errores
\item Corrige: $t = \lfloor (n-1)/2 \rfloor$ errores
\item \textbf{Ejemplo $(5,1)$:} Corrige hasta 2 errores de 5 bits
  \begin{itemize}
  \item Transmitimos 00000 o 11111
  \item Recibimos 01010 (2 errores) → mayoría 0 → decodificamos 0 ✓
  \item Recibimos 01110 (3 errores) → mayoría 1 → decodificamos 1 (error fatal ✗)
  \end{itemize}
\end{itemize}

\textbf{Estrategias de uso en la práctica:}

\textbf{Modo 1 - Solo detección (ARQ):}
\begin{itemize}
\item Usar toda la capacidad de detección: hasta $d_{min}-1$ errores
\item Cuando detectamos error → pedir retransmisión
\item \textbf{Ventaja:} Mayor débito útil (menos redundancia)
\item \textbf{Desventaja:} Retransmisiones aumentan latencia
\item \textbf{Uso:} Protocolos de red (TCP), almacenamiento (checksums)
\end{itemize}

\textbf{Modo 2 - Solo corrección (FEC):}
\begin{itemize}
\item Usar capacidad de corrección: hasta $\lfloor(d_{min}-1)/2\rfloor$ errores
\item Corregir automáticamente sin retransmisión
\item \textbf{Ventaja:} Sin retransmisiones, latencia fija
\item \textbf{Desventaja:} Más redundancia necesaria (menor $R_c$)
\item \textbf{Uso:} Broadcasting (TV, radio), tiempo real (VoIP), espacio
\end{itemize}

\textbf{Modo 3 - Híbrido (detección + corrección parcial):}
\begin{itemize}
\item Corregir $t$ errores, detectar $e > t$ errores adicionales
\item Requiere: $d_{min} \geq t + e + 1$
\item \textbf{Ejemplo con $d_{min}=7$:}
  \begin{itemize}
  \item Corregir 2 errores, detectar hasta 4 errores adicionales
  \item $t=2$, $e=4$ → $d_{min} \geq 2+4+1=7$ ✓
  \end{itemize}
\item \textbf{Ventaja:} Corrección automática para errores frecuentes, detección para casos raros
\item \textbf{Uso:} Sistemas modernos (LTE, 5G) con HARQ
\end{itemize}

\textbf{Trade-off fundamental (no se puede evitar):}
$$\text{Robusto (alto }d_{min}\text{)} \Leftrightarrow \text{Ineficiente (bajo }R_c\text{)}$$

\textbf{Cotas teóricas (límites matemáticos):}

\textbf{Cota de Singleton:} $d_{min} \leq n - k + 1$
\begin{itemize}
\item Límite superior en $d_{min}$ para tasa $R_c = k/n$ dada
\item Los códigos que alcanzan esta cota se llaman \textbf{MDS (Maximum Distance Separable)}
\item \textbf{Ejemplo:} Reed-Solomon codes (usados en CDs, DVDs, QR codes)
\end{itemize}

\textbf{Cota esfera-empaquetamiento (Hamming bound):}
$$2^k \sum_{i=0}^{t} \binom{n}{i} \leq 2^n$$

Límite en cuántas palabras código podemos "empaquetar" con esferas de radio $t$ sin solapamiento.

\textbf{Códigos perfectos:} Aquellos que alcanzan esta cota con igualdad.
\begin{itemize}
\item Códigos Hamming
\item Código Golay $(23,12)$
\item Código de repetición $(n,1)$ con $n$ impar
\end{itemize}

\textbf{Estos son muy raros - se conocen solo unos pocos!}

\subsection{Decodificación}

\textbf{Síndrome:}
$$\mathbf{s} = \mathbf{r}\mathbf{H}^T = (\mathbf{c} + \mathbf{e})\mathbf{H}^T = \mathbf{e}\mathbf{H}^T$$

\begin{itemize}
\item $\mathbf{s} = \mathbf{0}$: no errores detectados
\item $\mathbf{s} \neq \mathbf{0}$: error detectado
\end{itemize}

\textbf{Regla ML:} Elegir palabra código más cercana (mínima $d_H$).

\subsection{Códigos específicos}

\textbf{Repetición (n,1):}
\begin{itemize}
\item $d_{min} = n$, $t_c = \lfloor (n-1)/2 \rfloor$
\item Muy robusto, muy ineficiente ($R_c = 1/n$)
\end{itemize}

\textbf{Paridad (n,n-1):}
\begin{itemize}
\item $d_{min} = 2$, detecta 1 error, no corrige
\item Alta tasa ($R_c = (n-1)/n$)
\end{itemize}

\textbf{Hamming (2$^m$-1, 2$^m$-m-1):}
\begin{itemize}
\item $d_{min} = 3$, $t_c = 1$
\item Columnas de H: representación binaria de 1 a $2^m-1$
\item Ejemplo (7,4): $R_c = 4/7 \approx 0.57$
\end{itemize}

\subsection{Límites}

\textbf{Cota de Singleton:} $d_{min} \leq n - k + 1$

\textbf{Cota de Plotkin:} $d_{min} \leq \frac{n2^{k-1}}{2^k-1}$

\section{Canal Multitrayecto - La realidad física de la propagación}

\subsection{Motivación - ¿Por qué es complejo el canal real?}

\begin{tcolorbox}[colback=rojoclaro!10,colframe=rojoclaro,title=🧠 FUNDAMENTO FÍSICO]
\textbf{Esto era así porque:} En el mundo real, las ondas electromagnéticas NO viajan en línea recta desde el transmisor al receptor.

\textbf{Fenómenos físicos reales:}
\begin{enumerate}
\item \textbf{Reflexión:} La señal rebota en edificios, montañas, vehículos
\item \textbf{Difracción:} La señal rodea obstáculos (esquinas, colinas)
\item \textbf{Dispersión (scattering):} La señal se dispersa en objetos pequeños (árboles, farolas)
\item \textbf{Refracción:} La señal se dobla al atravesar capas atmosféricas diferentes
\end{enumerate}

\textbf{Consecuencia fundamental:} La señal llega al receptor por \textbf{múltiples caminos (multitrayecto)} con diferentes:
\begin{itemize}
\item \textbf{Retardos} $\tau_l$: cada camino tiene longitud diferente
\item \textbf{Atenuaciones} $\alpha_l$: cada camino pierde potencia distinto
\item \textbf{Fases} $\phi_l$: cada camino acumula desfase diferente
\end{itemize}

\textbf{Problema crítico:} Los caminos pueden interferirse \textbf{destructivamente}:
\begin{itemize}
\item Si dos caminos llegan en fase opuesta → se cancelan (desvanecimiento)
\item Si llegan desfasados en el tiempo → ISI (Interferencia Entre Símbolos)
\end{itemize}
\end{tcolorbox}

\textbf{Ejemplos prácticos donde multitrayecto es crítico:}
\begin{itemize}
\item \textbf{Comunicaciones móviles (4G/5G):} Señal rebota en edificios urbanos
  \begin{itemize}
  \item Delay spread: 0.1-20 $\mu$s (típico 1-5 $\mu$s)
  \item Puede haber 10-50 caminos significativos
  \end{itemize}
\item \textbf{WiFi indoor:} Reflexiones en paredes, muebles, personas
  \begin{itemize}
  \item Delay spread: 50-200 ns
  \item Canal muy variable (personas moviendose)
  \end{itemize}
\item \textbf{Comunicaciones submarinas acústicas:} Reflexiones en superficie y fondo marino
  \begin{itemize}
  \item Delay spread: decenas de ms (velocidad sonido = 1500 m/s)
  \item Canal \textbf{extremadamente} dispersivo
  \end{itemize}
\item \textbf{Broadcasting (TV, radio):} Reflexiones en terreno, ionosfera
  \begin{itemize}
  \item Delay spread: microsegundos a milisegundos
  \item Efecto "fantasma" en TV analógica
  \end{itemize}
\end{itemize}

\subsection{Modelo matemático del canal multitrayecto}

\textbf{Respuesta impulsional continua (física):}
$$\boxed{h_c(t;\tau) = \sum_{l=0}^{L-1} \alpha_l(t) e^{-i\phi_l(t)} \delta(\tau - \tau_l(t))}$$

Esta es la respuesta del canal en tiempo continuo, donde:
\begin{itemize}
\item $h_c(t;\tau)$: respuesta impulsional en tiempo $t$ con retardo $\tau$
\item $L$: número de caminos (trayectos) significativos
\item $\alpha_l(t)$: atenuación (ganancia) del camino $l$ en tiempo $t$
  \begin{itemize}
  \item Depende de $t$ porque objetos se mueven (Doppler)
  \item Típicamente: $\alpha_0 > \alpha_1 > \alpha_2 > \cdots$ (caminos más largos más atenuados)
  \end{itemize}
\item $\phi_l(t)$: fase del camino $l$ en tiempo $t$
  \begin{itemize}
  \item Incluye fase debido a longitud del camino: $\phi_l = 2\pi f_c \tau_l$
  \item Varía con $t$ por movimiento (efecto Doppler)
  \end{itemize}
\item $\tau_l(t)$: retardo de propagación del camino $l$
  \begin{itemize}
  \item $\tau_l = d_l / c$ donde $d_l$ es longitud del camino, $c$ velocidad de propagación
  \item Ordenados: $0 \leq \tau_0 < \tau_1 < \cdots < \tau_{L-1}$
  \end{itemize}
\item $\delta(\tau - \tau_l)$: Delta de Dirac (impulso en retardo $\tau_l$)
\end{itemize}

\textbf{Señal recibida (convolución con el canal):}
$$\boxed{r(t) = \int_{-\infty}^{+\infty} h_c(t;\tau) s(t-\tau) d\tau + n(t)}$$

Substituyendo la expresión de $h_c$:
$$\boxed{r(t) = \sum_{l=0}^{L-1} \alpha_l(t) e^{-i\phi_l(t)} s(t - \tau_l(t)) + n(t)}$$

\textbf{Interpretación física:}
\begin{itemize}
\item Recibimos una \textbf{suma ponderada} de copias retardadas de la señal transmitida
\item Cada copia tiene amplitud $\alpha_l$, fase $\phi_l$ y retardo $\tau_l$ diferentes
\item Las copias pueden sumarse \textbf{constructiva} o \textbf{destructivamente}
\item El ruido AWGN $n(t)$ se suma encima de toda esta mezcla
\end{itemize}

\textbf{Ejemplo numérico simple (2 caminos):}
\begin{itemize}
\item Camino directo: $\alpha_0 = 1.0$, $\tau_0 = 0$ (referencia)
\item Camino reflejado: $\alpha_1 = 0.5$, $\tau_1 = 1\,\mu$s
\item Señal recibida: $r(t) = 1.0 \cdot s(t) + 0.5 \cdot s(t - 1\,\mu\text{s}) + n(t)$
\item Si el símbolo dura $T_s = 0.5\,\mu$s → el eco llega 2 símbolos después → ISI severa
\end{itemize}

\subsection{Parámetros fundamentales del canal - Caracterización completa}

\begin{tcolorbox}[colback=azulclaro!10,colframe=azuloscuro,title=📏 Parámetros en Dominio TEMPORAL]
\textbf{1. Delay Spread (Dispersión temporal) $\tau_{max}$:}
$$\boxed{\tau_{max} = \tau_{L-1} - \tau_0}$$

Diferencia entre el camino más largo y el más corto.

\textbf{Esto era así porque:} Mide cuánto se "estira" la señal en el tiempo.

\textbf{Valores típicos:}
\begin{itemize}
\item Canal indoor (WiFi): 50-200 ns
\item Canal urbano (4G): 1-5 $\mu$s
\item Canal rural (cobertura extensa): 10-20 $\mu$s
\item Canal submarino: 10-100 ms (!)
\end{itemize}

\textbf{Impacto en ISI:}
\begin{itemize}
\item Si $\tau_{max} \ll T_s$ → ISI despreciable (canal "plano")
\item Si $\tau_{max} \approx T_s$ → ISI moderada (ecualizar simple)
\item Si $\tau_{max} \gg T_s$ → ISI severa (OFDM necesario)
\end{itemize}

\textbf{2. RMS Delay Spread (métrica estadística):}
$$\tau_{rms} = \sqrt{\frac{\sum_l \alpha_l^2 (\tau_l - \bar{\tau})^2}{\sum_l \alpha_l^2}}$$

donde $\bar{\tau} = \frac{\sum_l \alpha_l^2 \tau_l}{\sum_l \alpha_l^2}$ es retardo medio ponderado.

Mide la "dispersión promedio" considerando energía de cada camino.
\end{tcolorbox}

\begin{tcolorbox}[colback=verdeclaro!10,colframe=verdeclaro,title=📊 Parámetros en Dominio FRECUENCIAL]
\textbf{3. Ancho de banda de coherencia $B_c$:}
$$\boxed{B_c \approx \frac{1}{\tau_{max}}}$$

\textbf{Esto era así porque:} Es la transformada de Fourier del delay spread.

\textbf{Interpretación:} Ancho de banda en el cual el canal es aproximadamente \textbf{constante} ("plano").

\textbf{Relación con selectividad frecuencial:}
\begin{itemize}
\item Si $B_{se\tilde{n}al} \ll B_c$ → canal \textbf{plano en frecuencia} (flat fading)
  \begin{itemize}
  \item Todas las frecuencias afectadas igual
  \item No hay ISI significativa
  \item Ejemplo: GSM (200 kHz) en canal urbano ($B_c \approx$ 200 kHz)
  \end{itemize}
\item Si $B_{se\tilde{n}al} \gg B_c$ → canal \textbf{selectivo en frecuencia} (frequency-selective)
  \begin{itemize}
  \item Diferentes frecuencias atenuadas diferente
  \item ISI severa en tiempo
  \item Ejemplo: LTE (20 MHz) en canal urbano ($B_c \approx$ 200 kHz) → 100 veces mayor
  \end{itemize}
\end{itemize}

\textbf{Ejemplo numérico:}
\begin{itemize}
\item Delay spread: $\tau_{max} = 5\,\mu$s
\item Ancho coherencia: $B_c \approx 1/(5\mu\text{s}) = 200$ kHz
\item Si transmitimos con $B = 20$ MHz → $B/B_c = 100$ subcanales independientes
\item Por eso OFDM usa 100+ subportadoras (cada una dentro de $B_c$)
\end{itemize}
\end{tcolorbox}

\begin{tcolorbox}[colback=rojoclaro!10,colframe=rojoclaro,title=🚀 Parámetros por MOVILIDAD (Doppler)]
\textbf{4. Doppler Spread (Dispersión Doppler) $\nu_D$:}
$$\boxed{\nu_D = \frac{v_{max} \cdot f_c}{c}}$$

donde:
\begin{itemize}
\item $v_{max}$: velocidad relativa máxima transmisor-receptor [m/s]
\item $f_c$: frecuencia portadora [Hz]
\item $c$: velocidad de propagación [m/s]
  \begin{itemize}
  \item Ondas radio: $c = 3 \times 10^8$ m/s
  \item Sonido submarino: $c \approx 1500$ m/s
  \end{itemize}
\end{itemize}

\textbf{Esto era así porque:} Efecto Doppler causa desplazamiento de frecuencia proporcional a velocidad.

\textbf{Ejemplos prácticos:}
\begin{itemize}
\item \textbf{Peatón (3 km/h = 0.83 m/s), $f_c = 2$ GHz:}
  $$\nu_D = \frac{0.83 \times 2 \times 10^9}{3 \times 10^8} \approx 5.5 \text{ Hz}$$
\item \textbf{Coche ciudad (50 km/h = 13.9 m/s), $f_c = 2$ GHz:}
  $$\nu_D \approx 93 \text{ Hz}$$
\item \textbf{Tren alta velocidad (300 km/h = 83.3 m/s), $f_c = 2$ GHz:}
  $$\nu_D \approx 555 \text{ Hz}$$
\item \textbf{Submarino (20 nudos $\approx$ 10 m/s), $f_c = 10$ kHz, $c = 1500$ m/s:}
  $$\nu_D = \frac{10 \times 10^4}{1500} \approx 67 \text{ Hz}$$ (!)
\end{itemize}

\textbf{5. Tiempo de coherencia $T_c$ (cuánto dura la "estacionariedad"):}
$$\boxed{T_c \approx \frac{1}{\nu_D}}$$

\textbf{Interpretación:} Tiempo durante el cual el canal permanece aproximadamente \textbf{constante}.

\textbf{Relación con variabilidad temporal:}
\begin{itemize}
\item Si $T_{símbolo} \ll T_c$ → canal \textbf{cuasi-estático} (slow fading)
  \begin{itemize}
  \item Canal constante durante muchos símbolos
  \item Puedo estimar canal y usar estimación durante tiempo
  \item Ejemplo: WiFi con usuario fijo
  \end{itemize}
\item Si $T_{símbolo} \gg T_c$ → canal \textbf{rápidamente variable} (fast fading)
  \begin{itemize}
  \item Canal cambia dentro de un símbolo
  \item Necesito re-estimar constantemente
  \item Ejemplo: 4G en tren alta velocidad
  \end{itemize}
\end{itemize}

\textbf{Ejemplos numéricos:}
\begin{itemize}
\item Peatón ($\nu_D = 5.5$ Hz): $T_c \approx 182$ ms (muy largo)
\item Coche ($\nu_D = 93$ Hz): $T_c \approx 11$ ms
\item Tren ($\nu_D = 555$ Hz): $T_c \approx 1.8$ ms (corto!)
\end{itemize}

\textbf{Relación con OFDM:}
\begin{itemize}
\item Símbolo OFDM típico: $T_{OFDM} \approx 50$-$100\,\mu$s
\item Peatón: $T_c/T_{OFDM} \approx 1820$ símbolos (super estable)
\item Tren: $T_c/T_{OFDM} \approx 18$ símbolos (necesita pilotos frecuentes)
\end{itemize}
\end{tcolorbox}

\subsection{Canal discreto equivalente - De continuo a discreto}

\begin{tcolorbox}[colback=azulclaro!10,colframe=azuloscuro,title=🧠 FUNDAMENTO: ¿Por qué "discreto"?]
\textbf{Esto era así porque:} Después de filtrado y muestreo, trabajamos con \textbf{secuencias discretas} (no señales continuas).

\textbf{Cadena completa:}
\begin{enumerate}
\item Transmisor emite señal continua: $s(t) = \sum_k a_k h_e(t - kT_s)$
\item Canal multitrayecto: $r(t) = \sum_l \alpha_l s(t - \tau_l) + n(t)$
\item Receptor aplica filtro adaptado: $y(t) = r(t) * g_r(t)$
\item Muestreo en instantes óptimos: $y_k = y(t_0 + kT_s)$
\end{enumerate}

\textbf{Resultado:} Obtenemos un \textbf{modelo discreto} que relaciona símbolos transmitidos $\{a_k\}$ con muestras recibidas $\{y_k\}$.
\end{tcolorbox}

\textbf{Modelo del canal discreto equivalente:}
$$\boxed{y_k = \sum_{l=0}^{L-1} h_l a_{k-l} + b_k}$$

donde:
\begin{itemize}
\item $y_k$: muestra recibida en instante $k$ (después de filtrado y muestreo)
\item $h_l$: coeficiente $l$ del canal discreto (complejo en general)
  \begin{itemize}
  \item Captura efecto combinado de: pulso transmisor + canal físico + filtro receptor
  \item $h_0$ es coeficiente del camino "principal" (típicamente el más fuerte)
  \item $|h_l|$ decrece generalmente con $l$ (caminos largos más atenuados)
  \end{itemize}
\item $a_{k-l}$: símbolo transmitido en instante $k-l$ (de la constelación)
\item $L$: longitud de memoria del canal (número de "taps" o coeficientes)
  \begin{itemize}
  \item Relacionado con delay spread: $L = \lceil \tau_{max}/T_s \rceil + 1$
  \item Si $\tau_{max} = 5\,\mu$s y $T_s = 1\,\mu$s → $L = 6$ taps
  \end{itemize}
\item $b_k$: ruido gaussiano blanco discreto, $b_k \sim \mathcal{CN}(0, N_0)$
  \begin{itemize}
  \item Proviene del ruido AWGN después de filtrado
  \item Sigue siendo gaussiano (filtro lineal preserva gaussianidad)
  \item Puede estar coloreado si no usamos matched filter
  \end{itemize}
\end{itemize}

\textbf{Cómo se calculan los coeficientes $h_l$:}

Definimos la respuesta global del sistema:
$$r(t) = (h_e * c * g_r)(t)$$

donde $c(t)$ es la respuesta del canal físico, $h_e(t)$ es el pulso transmisor, $g_r(t)$ es el filtro receptor.

\textbf{Los coeficientes discretos son:}
$$\boxed{h_l = r(t_0 + lT_s)}$$

es decir, muestras de la respuesta global en los instantes de símbolo.

\textbf{Ejemplo práctico (matched filter con 2 caminos):}

Supongamos:
\begin{itemize}
\item Pulso rectangular: $h_e(t) = \text{rect}(t/T)$
\item Canal: $c(t) = \delta(t) + 0.5\delta(t - 1.5T)$
\item Matched filter: $g_r(t) = h_e(-t)$
\item Respuesta global: $r(t) = \Lambda(t/T) + 0.5\Lambda((t-1.5T)/T)$ (triangular)
\end{itemize}

Muestreando en $t = kT$:
\begin{align*}
h_0 &= r(0) = 1.0\\
h_1 &= r(T) = 0.25 \quad \text{(contribución del eco)}\\
h_2 &= r(2T) = 0.25\\
h_l &= 0 \quad \text{para } l \geq 3
\end{align*}

Por tanto: $L = 3$ (canal con 3 taps).

\textbf{Representación en transformada Z:}
$$H(z) = \sum_{l=0}^{L-1} h_l z^{-l}$$

Esto es un \textbf{filtro FIR} (Finite Impulse Response) que modela el canal.

\textbf{Propiedades importantes:}
\begin{itemize}
\item Si $L = 1$ (solo $h_0 \neq 0$): canal sin memoria, no hay ISI
\item Si $L > 1$: canal con memoria, hay ISI (necesitamos ecualizar)
\item Los coeficientes $h_l$ varían lentamente con tiempo (por movilidad)
  \begin{itemize}
  \item Necesitamos re-estimar cada $T_c$ (tiempo de coherencia)
  \end{itemize}
\end{itemize}

\textbf{Estimación de los coeficientes del canal:}

\textbf{Método 1: Correlación cruzada (secuencia de entrenamiento):}
$$R_{ya}(u) = \sum_k y_k a^*(k-u)$$

donde $\{a_k\}$ es una secuencia de entrenamiento \textbf{conocida} por el receptor.

\textbf{Relación fundamental:}
$$\boxed{R_{ya}(u) = \sum_{l=0}^{L-1} h_l R_{aa}(u-l) = (h * R_{aa})(u)}$$

donde $R_{aa}(u) = \sum_k a_k a^*_{k-u}$ es la autocorrelación de la secuencia.

\textbf{Caso ideal:} Si usamos una secuencia con $R_{aa}(u) = \delta(u)$ (impulso discreto):
$$\boxed{\hat{h}_u = R_{ya}(u) \quad \text{(estimación directa!)}}$$

\textbf{Secuencias que cumplen esto aproximadamente:}
\begin{itemize}
\item Secuencias PN (Pseudo-Noise): ruido pseudoaleatorio
\item Secuencias de Zadoff-Chu (usadas en LTE)
\item Secuencias de Gold (usadas en GPS)
\end{itemize}

\textbf{Método 2: Solución de mínimos cuadrados (LS):}

Formulamos sistema lineal:
$$\mathbf{y} = \mathbf{A}\mathbf{h} + \mathbf{n}$$

donde $\mathbf{A}$ es matriz de Toeplitz con la secuencia de entrenamiento.

\textbf{Solución óptima LS:}
$$\boxed{\hat{\mathbf{h}} = (\mathbf{A}^H \mathbf{A})^{-1} \mathbf{A}^H \mathbf{y}}$$

Esta solución minimiza el error cuadrático medio.

\section{Ecualización - Combatir la ISI}

\begin{tcolorbox}[colback=rojoclaro!10,colframe=rojoclaro,title=🧠 MOTIVACIÓN PROFUNDA]
\textbf{Problema fundamental:} El canal multitrayecto causa ISI (Interferencia Entre Símbolos).

\textbf{Esto era así porque:} Cada símbolo recibido contiene contribuciones de símbolos pasados (y futuros):
$$y_k = \underbrace{h_0 a_k}_{\text{Señal deseada}} + \underbrace{\sum_{l\neq 0} h_l a_{k-l}}_{\text{ISI (indeseado)}} + b_k$$

\textbf{Objetivo de la ecualización:} Diseñar un filtro $g(t)$ en el receptor para:
\begin{enumerate}
\item \textbf{Cancelar la ISI:} Hacer que la respuesta global cumpla Nyquist
\item \textbf{No amplificar mucho el ruido:} Mantener SNR razonable
\end{enumerate}

\textbf{Trade-off inevitable:}
$$\text{Cancelar ISI} \Leftrightarrow \text{Amplificar ruido}$$

Diferentes estrategias de ecualización priorizan una u otra.
\end{tcolorbox}

\subsection{Zero-Forcing (ZF) - Eliminación perfecta de ISI}

\textbf{Idea central:} Diseñar un ecualizador que sea la \textbf{inversa exacta} del canal.

\textbf{Criterio matemático:}
$$G(z) = \frac{1}{H(z)} \quad \Leftrightarrow \quad g * h = \delta$$

donde $H(z)$ es la transformada Z del canal, $G(z)$ es la del ecualizador.

\textbf{Respuesta global del sistema ecualizado:}
$$(h * g)(k) = \delta(k) = \begin{cases} 1 & k=0\\ 0 & k \neq 0 \end{cases}$$

Perfecto! No hay ISI en absoluto (Nyquist criterion satisfied).

\textbf{Implementación práctica:}

Si $H(z) = h_0 + h_1 z^{-1} + \cdots + h_{L-1}z^{-(L-1)}$, entonces:
$$G(z) = \frac{1}{h_0(1 + (h_1/h_0)z^{-1} + \cdots)}$$

Usando desarrollo en serie (si converge):
$$G(z) = \frac{1}{h_0}(1 - (h_1/h_0)z^{-1} + ((h_1/h_0)^2 - h_2/h_0)z^{-2} + \cdots)$$

Puede requerir un filtro IIR (Infinite Impulse Response) de longitud infinita.

\textbf{\textcolor{red}{Problema crítico - Amplificación de ruido:}}

Supongamos que el canal tiene un \textbf{cero (nulo) en frecuencia} $\omega_0$:
$$H(e^{i\omega_0}) \approx 0$$

Entonces el ecualizador ZF tendrá:
$$G(e^{i\omega_0}) = \frac{1}{H(e^{i\omega_0})} \to \infty$$

¡El ecualizador \textbf{amplifica infinitamente} el ruido en esa frecuencia!

\textbf{Ejemplo numérico:}
\begin{itemize}
\item Canal: $H(z) = 1 + 0.9z^{-1}$ (tiene cero cerca de $z=-1$)
\item Ecualizador ZF: $G(z) = \frac{1}{1+0.9z^{-1}} = 1 - 0.9z^{-1} + 0.81z^{-2} - \cdots$
\item En frecuencia $\omega = \pi$: $|G(e^{i\pi})| = \frac{1}{|1-0.9|} = 10$ (amplifica $\times$10!)
\item El ruido en esa banda se multiplica por 10 → SNR cae drásticamente
\end{itemize}

\textbf{Ventajas de ZF:}
\begin{itemize}
\item ISI = 0 (eliminación perfecta)
\item Simple conceptualmente
\item No necesita conocer estadísticas del ruido
\end{itemize}

\textbf{Desventajas de ZF:}
\begin{itemize}
\item \textbf{Amplificación severa de ruido} (especialmente si canal tiene nulos profundos)
\item Puede ser inestable si $H(z)$ tiene ceros fuera del círculo unitario
\item SNR de salida puede ser muy bajo
\item Rendimiento pobre en canales malos
\end{itemize}

\textbf{Cuándo usar ZF:}
\begin{itemize}
\item Canal "bueno" sin nulos profundos
\item SNR de entrada muy alto (ruido ya despreciable)
\item Cuando ISI es el problema dominante
\end{itemize}

\subsection{MMSE (Minimum Mean Square Error) - Compromiso óptimo}

\textbf{Idea central:} Minimizar el \textbf{error cuadrático medio} entre símbolo transmitido y estimado.

\textbf{Criterio de optimización:}
$$\min_{\mathbf{g}} \text{MSE} = \min_{\mathbf{g}} E\left[|a_k - \hat{a}_{k-\Delta}|^2\right]$$

donde $\hat{a}_{k-\Delta} = \mathbf{g}^H \mathbf{y}_k$ es la estimación del símbolo con retardo $\Delta$.

\textbf{Solución óptima (ecuaciones de Wiener-Hopf):}
$$\boxed{\mathbf{g}_{MMSE} = (\sigma_a^2 \mathbf{H}^H\mathbf{H} + \sigma_n^2\mathbf{I})^{-1}\mathbf{H}^H\mathbf{e}_{\Delta}}$$

donde:
\begin{itemize}
\item $\mathbf{H}$: matriz de canal (convolución como multiplicación matricial)
\item $\sigma_a^2$: potencia promedio de los símbolos transmitidos
\item $\sigma_n^2$: potencia del ruido
\item $\mathbf{e}_{\Delta}$: vector canónico (0,\ldots,0,1,0,\ldots,0) en posición $\Delta$
\end{itemize}

\textbf{Interpretación intuitiva:}
$$\mathbf{g}_{MMSE} = \underbrace{(\mathbf{H}^H\mathbf{H} + \frac{\sigma_n^2}{\sigma_a^2}\mathbf{I})^{-1}\mathbf{H}^H}_{\text{Regularización}} \mathbf{e}_{\Delta}$$

La matriz $(\mathbf{H}^H\mathbf{H} + \lambda\mathbf{I})^{-1}$ con $\lambda = \sigma_n^2/\sigma_a^2$ es una \textbf{regularización}:
\begin{itemize}
\item Cuando SNR alto ($\sigma_n^2 \to 0$): $\mathbf{g}_{MMSE} \to \mathbf{g}_{ZF}$ (comportamiento ZF)
\item Cuando SNR bajo ($\sigma_n^2$ grande): $\mathbf{g}_{MMSE}$ evita amplificar ruido excesivamente
\end{itemize}

\textbf{En dominio frecuencial (forma simple):}
$$G_{MMSE}(\omega) = \frac{H^*(\omega)}{|H(\omega)|^2 + \sigma_n^2/\sigma_a^2}$$

Compara con ZF: $G_{ZF}(\omega) = 1/H(\omega)$

\textbf{Diferencia clave:}
\begin{itemize}
\item ZF invierte directamente $H$ (puede explotar si $H \approx 0$)
\item MMSE invierte $H$ pero añade $\sigma_n^2/\sigma_a^2$ en denominador (estabiliza)
\end{itemize}

\textbf{Ventajas de MMSE:}
\begin{itemize}
\item \textbf{Compromiso óptimo} entre ISI y amplificación de ruido (minimiza MSE)
\item Siempre estable (no invierte directamente el canal)
\item Mejor rendimiento que ZF en SNR bajo/medio
\item Se adapta automáticamente al SNR
\end{itemize}

\textbf{Desventajas de MMSE:}
\begin{itemize}
\item ISI residual (no elimina completamente ISI)
\item Necesita conocer $\sigma_n^2$ y $\sigma_a^2$ (estadísticas)
\item Complejidad computacional mayor (inversión matricial)
\end{itemize}

\textbf{Cuándo usar MMSE:}
\begin{itemize}
\item Escenarios generales (es el estándar en sistemas modernos)
\item SNR bajo/medio (mejor que ZF)
\item Canal con nulos profundos (MMSE no explota)
\end{itemize}

\subsection{Decision Feedback Equalizer (DFE) - Uso inteligente de decisiones}

\textbf{Idea brillante:} Usar decisiones \textbf{pasadas} (ya tomadas) para cancelar ISI de forma no lineal.

\textbf{Estructura en dos partes:}
\begin{enumerate}
\item \textbf{Filtro Feedforward $g_f$:} Procesa señal recibida $\mathbf{y}$
\item \textbf{Filtro Feedback $g_b$:} Usa decisiones pasadas $\hat{\mathbf{a}}$
\end{enumerate}

\textbf{Ecuación del DFE:}
$$\boxed{z_k = \underbrace{\sum_i g_{f,i} y_{k-i}}_{\text{Feedforward}} - \underbrace{\sum_{j=1}^{L_b} g_{b,j} \hat{a}_{k-j}}_{\text{Feedback}}}$$

Luego: $\hat{a}_k = \text{decision}(z_k)$ (umbral o más cercano en constelación)

\textbf{¿Por qué funciona esto?}

\begin{tcolorbox}[colback=verdeclaro!10,colframe=verdeclaro,title=💡 Intuición del DFE]
\textbf{Recordatorio:} La ISI viene de símbolos pasados y futuros:
$$y_k = h_0 a_k + \underbrace{\sum_{l<0} h_l a_{k-l}}_{\text{Pre-cursor ISI}} + \underbrace{\sum_{l>0} h_l a_{k-l}}_{\text{Post-cursor ISI}} + b_k$$

\textbf{Estrategia del DFE:}
\begin{itemize}
\item \textbf{Feedforward:} Cancela pre-cursor ISI (símbolos futuros) + ajusta fase/amplitud
\item \textbf{Feedback:} Cancela post-cursor ISI (símbolos pasados) usando $\hat{a}_{k-1}, \hat{a}_{k-2}, \ldots$ (ya decididos)
\end{itemize}

\textbf{Ventaja clave del feedback:}
\begin{itemize}
\item Las decisiones $\hat{a}_k \in \{$constelación$\}$ son \textbf{deterministas} (no ruido)
\item Por tanto, el filtro feedback \textbf{NO amplifica ruido} (solo procesa decis ones binarias/discretas)
\item ZF/MMSE amplifican ruido al invertir canal, DFE no tiene este problema en parte feedback
\end{itemize}
\end{tcolorbox}

\textbf{Diseño óptimo (criterio MMSE):}
Minimizar MSE = $E[|a_k - z_k|^2]$ eligiendo $g_f$ y $g_b$ óptimos.

\textbf{Solución:} Ecuaciones de Wiener-Hopf no lineales (más complejas que MMSE lineal).

\textbf{Ventajas del DFE:}
\begin{itemize}
\item \textbf{Mejor que ecualización lineal} (ZF/MMSE) en canales severos
\item No amplifica ruido en parte feedback
\item Puede lograr casi ZF performance sin penalty de ruido
\item Rendimiento cercano a óptimo (ML/Viterbi) con menor complejidad
\end{itemize}

\textbf{\textcolor{red}{Problema crítico - Propagación de errores:}}

Si una decisión $\hat{a}_{k-1}$ es \textbf{errónea}:
\begin{enumerate}
\item Feedback cancela ISI basado en $\hat{a}_{k-1}$ incorrecto
\item Esto corrompe $z_k$ → decisión $\hat{a}_k$ también errónea (alta probabilidad)
\item Que corrompe $z_{k+1}$ → decisión $\hat{a}_{k+1}$ errónea
\item ... \textbf{error burst} (ráfaga de errores)
\end{enumerate}

\textbf{Ejemplo:} Un solo error puede causar 5-10 errores consecutivos.

\textbf{Soluciones a la propagación:}
\begin{enumerate}
\item \textbf{Secuencias de entrenamiento periódicas:} "Resetear" el feedback cada cierto tiempo
\item \textbf{Códigos correctores:} FEC puede corregir los burst errors
\item \textbf{Interleaving:} Dispersar errores en el tiempo
\item \textbf{Usar Viterbi en su lugar:} No propaga errores (pero más complejo)
\end{enumerate}

\textbf{Cuándo usar DFE:}
\begin{itemize}
\item Canales con ISI severa (mejor que MMSE lineal)
\item Cuando complejidad Viterbi es prohibitiva
\item Con código corrector para manejar error propagation
\item DSL, cable modems, comunicaciones por cable
\end{itemize}

\subsection{Maximum Likelihood / Viterbi - Solución óptima}

\textbf{Idea fundamental:} En vez de decidir símbolo por símbolo, buscar la \textbf{secuencia completa más probable}.

\textbf{Criterio ML:}
$$\hat{\mathbf{a}} = \arg\max_{\mathbf{a}} p(\mathbf{y} | \mathbf{a}, \mathbf{H})$$

Buscar entre \textbf{todas} las posibles secuencias la que maximiza verosimilitud.

\textbf{Problema:} Si transmitimos $N$ símbolos con alfabeto $M$ → hay $M^N$ secuencias posibles.
\begin{itemize}
\item Para QPSK ($M=4$), $N=100$ → $4^{100} \approx 10^{60}$ secuencias (!)
\item Búsqueda exhaustiva: \textbf{imposible}
\end{itemize}

\textbf{Solución: Algoritmo de Viterbi (programación dinámica)}

Explota la estructura del canal con memoria finita mediante un \textbf{trellis}.

\textbf{Trellis (enrejado):}
\begin{itemize}
\item \textbf{Estados:} Memoria del canal $S_k = (a_{k-1}, a_{k-2}, \ldots, a_{k-L+1})$
\item Número de estados: $M^{L-1}$
\item \textbf{Transiciones:} Símbolo actual $a_k$ causa transición $S_k \to S_{k+1}$
\item \textbf{Métrica de rama:} Distancia euclidiana $\|y_k - \hat{y}_k(S_k, a_k)\|^2$
\end{itemize}

\textbf{Ejemplo (canal con 1 tap de memoria, BPSK):}
\begin{itemize}
\item Modelo: $y_k = h_0 a_k + h_1 a_{k-1} + b_k$ con $a_k \in \{-1,+1\}$
\item Estados: $S_k = a_{k-1} \in \{-1, +1\}$ → 2 estados
\item Transiciones: De cada estado salen 2 ramas ($a_k = \pm 1$)
\item Complejidad: $O(2 \cdot 2 \cdot N) = O(4N)$ (lineal en $N$!)
\end{itemize}

\textbf{Complejidad general de Viterbi:}
$$\text{Complejidad} = O(M^L \cdot N)$$

donde $M^L$ es el número de estados.

\textbf{Problema:} Crece \textbf{exponencialmente} con la memoria del canal $L$.
\begin{itemize}
\item $L=2$, $M=4$ (QPSK): $4^2 = 16$ estados (manejable)
\item $L=5$, $M=4$: $4^5 = 1024$ estados (pesado)
\item $L=10$, $M=4$: $4^{10} \approx 10^6$ estados (prohibitivo)
\end{itemize}

\textbf{Ventajas de Viterbi/ML:}
\begin{itemize}
\item \textbf{Óptimo} (minimiza probabilidad de error de secuencia)
\item No propaga errores (considera secuencia completa)
\item Mejor rendimiento posible para canal conocido
\end{itemize}

\textbf{Desventajas:}
\begin{itemize}
\item Complejidad exponencial en $L$ (intratable para $L$ grande)
\item Requiere conocimiento perfecto del canal $\mathbf{H}$
\item Latencia (necesita observar ventana de símbolos)
\end{itemize}

\textbf{Cuándo usar Viterbi:}
\begin{itemize}
\item Canal con memoria corta ($L \leq 3$-$5$)
\item Alfabeto pequeño ($M \leq 4$)
\item Cuando performance es crítico (espacial, militar)
\item Decodificación de códigos convolucionales (uso principal)
\end{itemize}

\textbf{Comparación de ecualizadores (resumen):}

\begin{center}
\begin{tabular}{|l|c|c|c|c|}
\hline
\textbf{Método} & \textbf{ISI} & \textbf{Ruido} & \textbf{Complejidad} & \textbf{Errores} \\
\hline
ZF & Nula & Amplifica & Baja & No propaga \\
MMSE & Residual & Compromiso & Media & No propaga \\
DFE & Casi nula & No ampli. FB & Media & Propaga \\
Viterbi & Nula & Óptimo & Exponencial & No propaga \\
\hline
\end{tabular}
\end{center}

\textbf{Elección práctica:}
\begin{itemize}
\item SNR alto, $L$ pequeño: ZF (simple)
\item Caso general: MMSE o DFE (compromiso)
\item Performance crítico, $L$ pequeño: Viterbi (óptimo)
\item $L$ grande: OFDM (evita ecualización compleja)
\end{itemize}

\section{OFDM - Solución elegante al multipaso}

\subsection{Principio y motivación}

\textbf{OFDM (Orthogonal Frequency-Division Multiplexing)} es una técnica de transmisión multiportadora que divide el canal en muchas subportadoras ortogonales estrechas.

\textbf{¿Por qué OFDM?}

\textbf{Problema en canales multitrayecto:}
\begin{itemize}
\item La señal llega por múltiples caminos con diferentes retardos
\item Causa interferencia entre símbolos (ISI) severa
\item Ecualización en tiempo es muy compleja (filtros largos, adaptación difícil)
\end{itemize}

\textbf{Idea brillante de OFDM:}
\begin{itemize}
\item Dividir el canal ancho en $N$ subcanales estrechos
\item Cada subcanal es suficientemente estrecho para ser "plano" (sin ISI significativa)
\item Transmitir en paralelo sobre todas las subportadoras
\item Usar DFT/IDFT para modular/demodular (implementación eficiente con FFT)
\end{itemize}

\textbf{Analogía:} En vez de conducir un camión grande por una carretera con baches (difícil), enviamos muchos coches pequeños por carriles separados (fácil).

\textbf{Ventajas principales:}
\begin{itemize}
\item \textbf{Ecualización simple:} Un coeficiente complejo por subportadora (escalar, no vectorial)
\item \textbf{Robustez:} Si una subportadora sufre desvanecimiento profundo, solo afecta a parte de los datos
\item \textbf{Adaptación:} Podemos usar diferentes modulaciones en cada subportadora (bit loading)
\item \textbf{Flexibilidad:} Fácil asignar subportadoras a diferentes usuarios (OFDMA)
\end{itemize}

\textbf{Inconvenientes:}
\begin{itemize}
\item \textbf{PAPR elevado:} La suma de muchas portadoras crea picos de amplitud
\item \textbf{Sensibilidad a offset de frecuencia:} Destruye ortogonalidad (ICI - Inter-Carrier Interference)
\item \textbf{Overhead del CP:} Prefijo cíclico reduce eficiencia espectral
\end{itemize}

\textbf{Aplicaciones:} WiFi (802.11a/g/n/ac/ax), 4G/5G (LTE, NR), DVB-T (TV digital), DSL

\subsection{Modulación y demodulación - Magia de la DFT}

\begin{tcolorbox}[colback=rojoclaro!10,colframe=rojoclaro,title=🧠 FUNDAMENTO MATEMÁTICO]
\textbf{Esto era así porque:} OFDM explota la propiedad de \textbf{ortogonalidad} de exponenciales complejas.

\textbf{Recordatorio - Ortogonalidad continua:}
En el dominio continuo, las sinusoides $e^{i2\pi f_k t}$ con frecuencias $f_k = k/T$ son ortogonales en el intervalo $[0,T]$:
$$\int_0^T e^{i2\pi f_k t} e^{-i2\pi f_l t} dt = T\delta_{k,l}$$

\textbf{Versión discreta (base de OFDM):}
Las exponenciales discretas $e^{i2\pi kn/N}$ son ortogonales sobre $n \in \{0,1,\ldots,N-1\}$:
$$\boxed{\sum_{n=0}^{N-1} e^{i2\pi kn/N} e^{-i2\pi ln/N} = N\delta_{k,l}}$$

donde $\delta_{k,l} = \begin{cases} 1 & k=l\\ 0 & k\neq l \end{cases}$

\textbf{¿Por qué es crucial esto?} Permite que $N$ señales compartan el mismo espectro sin interferirse, multiplicando la capacidad por $N$.
\end{tcolorbox}

OFDM usa la DFT (Discrete Fourier Transform) para implementar eficientemente la modulación multiportadora.

\textbf{IDFT - Modulación en el transmisor (síntesis):}
$$\boxed{x_n = \frac{1}{\sqrt{N}}\sum_{k=0}^{N-1} s_k e^{i2\pi kn/N}}, \quad n = 0,\ldots,N-1$$

\textbf{Interpretación física paso a paso:}
\begin{enumerate}
\item \textbf{Entrada:} $N$ símbolos complejos $\{s_0, s_1, \ldots, s_{N-1}\}$
  \begin{itemize}
  \item Cada $s_k$ es un símbolo QAM/PSK en la subportadora $k$
  \item Por ejemplo, con QPSK: $s_k \in \{1+i, 1-i, -1+i, -1-i\}/\sqrt{2}$
  \end{itemize}
\item \textbf{Modulación:} Cada símbolo $s_k$ se multiplica por su "base" $e^{i2\pi kn/N}$
  \begin{itemize}
  \item Esta exponencial representa una subportadora de frecuencia $f_k = k/(NT_s)$
  \item En tiempo continuo equivalente: $e^{i2\pi f_k t}$ con $f_k = k\Delta f$
  \end{itemize}
\item \textbf{Suma coherente:} Se suman todas las contribuciones de las $N$ subportadoras
\item \textbf{Salida:} $N$ muestras temporales $\{x_0, x_1, \ldots, x_{N-1}\}$
  \begin{itemize}
  \item Estas son las muestras del símbolo OFDM en dominio tiempo
  \item Se transmiten secuencialmente a tasa $1/T_s$
  \end{itemize}
\end{enumerate}

\textbf{Ejemplo numérico simple ($N=4$):}

Supongamos $s = [1, i, -1, -i]$ (4 símbolos QPSK en 4 subportadoras):
\begin{align*}
x_0 &= \frac{1}{2}(1 \cdot e^0 + i \cdot e^0 + (-1) \cdot e^0 + (-i) \cdot e^0) = 0\\
x_1 &= \frac{1}{2}(1 \cdot e^{i0} + i \cdot e^{i\pi/2} + (-1) \cdot e^{i\pi} + (-i) \cdot e^{i3\pi/2}) = 1\\
x_2 &= \frac{1}{2}(1 + i \cdot e^{i\pi} + (-1) \cdot e^{i2\pi} + (-i) \cdot e^{i3\pi}) = 0\\
x_3 &= \frac{1}{2}(1 + i \cdot e^{i3\pi/2} + (-1) \cdot e^{i3\pi} + (-i) \cdot e^{i9\pi/2}) = -1
\end{align*}

Símbolo OFDM temporal: $x = [0, 1, 0, -1]$

\textbf{DFT - Demodulación en el receptor (análisis):}
$$\boxed{\hat{s}_k = \frac{1}{\sqrt{N}}\sum_{n=0}^{N-1} y_n e^{-i2\pi kn/N}}, \quad k = 0,\ldots,N-1$$

\textbf{Interpretación física:}
\begin{enumerate}
\item \textbf{Entrada:} $N$ muestras recibidas $\{y_0, y_1, \ldots, y_{N-1}\}$
  \begin{itemize}
  \item Después de filtrado, muestreo y remoción del CP
  \end{itemize}
\item \textbf{Proyección:} Cada muestra $y_n$ se multiplica por $e^{-i2\pi kn/N}$ (conjugado de la base)
  \begin{itemize}
  \item Esto "correlaciona" la señal recibida con cada subportadora
  \item Similar a un matched filter para cada frecuencia
  \end{itemize}
\item \textbf{Integración discreta:} Se suman las $N$ contribuciones
  \begin{itemize}
  \item Por ortogonalidad, solo sobrevive la componente de frecuencia $k$
  \item Las otras subportadoras se cancelan (suma a cero)
  \end{itemize}
\item \textbf{Salida:} $N$ símbolos estimados $\{\hat{s}_0, \hat{s}_1, \ldots, \hat{s}_{N-1}\}$
  \begin{itemize}
  \item Cada $\hat{s}_k$ es la estimación del símbolo en subportadora $k$
  \item Luego se demodula cada uno según constelación (QAM/PSK)
  \end{itemize}
\end{enumerate}

\textbf{Verificación de la ortogonalidad (demo matemática):}

Si transmitimos símbolo $s_k$ en subportadora $k$, recibimos:
\begin{align*}
y_n &= x_n = \frac{1}{\sqrt{N}} s_k e^{i2\pi kn/N} \quad \text{(sin canal, sin ruido)}\\
\hat{s}_l &= \frac{1}{\sqrt{N}}\sum_{n=0}^{N-1} y_n e^{-i2\pi ln/N}\\
&= \frac{1}{N} s_k \sum_{n=0}^{N-1} e^{i2\pi (k-l)n/N}\\
&= \frac{1}{N} s_k \cdot N\delta_{k,l} = s_k \delta_{k,l}
\end{align*}

Por tanto: $\hat{s}_l = \begin{cases} s_k & l=k\\ 0 & l \neq k \end{cases}$

¡Perfecto! Solo recuperamos el símbolo de la subportadora deseada, sin interferencia.

\textbf{Implementación práctica - FFT (Fast Fourier Transform):}

\begin{tcolorbox}[colback=verdeclaro!10,colframe=verdeclaro,title=⚡ Eficiencia Computacional]
\textbf{DFT directa:} $O(N^2)$ operaciones (multiplicaciones complejas)
\begin{itemize}
\item Para cada $k$ (hay $N$ valores): sumamos $N$ términos
\item Total: $N \times N = N^2$ multiplicaciones
\end{itemize}

\textbf{FFT (algoritmo Cooley-Tukey):} $O(N\log_2 N)$ operaciones
\begin{itemize}
\item Explota simetría y periodicidad de $e^{i2\pi/N}$
\item Divide problema recursivamente (divide-and-conquer)
\item Funciona mejor cuando $N = 2^m$ (potencia de 2)
\end{itemize}

\textbf{Ganancia práctica:}
\begin{itemize}
\item $N=64$: DFT = 4096 ops, FFT = 384 ops → $\times 10.7$ más rápido
\item $N=256$: DFT = 65536 ops, FFT = 2048 ops → $\times 32$ más rápido
\item $N=1024$: DFT = 1048576 ops, FFT = 10240 ops → $\times 102$ más rápido
\item $N=2048$: DFT = 4194304 ops, FFT = 22528 ops → $\times 186$ más rápido
\end{itemize}

\textbf{Sin FFT, OFDM sería impracticable en tiempo real!}
\end{tcolorbox}

\textbf{Valores típicos de $N$ en estándares:}
\begin{itemize}
\item \textbf{WiFi 802.11a/g:} $N = 64$ (52 datos + 4 pilotos + 8 nulos)
\item \textbf{WiFi 802.11n:} $N = 64$ (40 MHz) o $N = 128$ (80 MHz)
\item \textbf{LTE:} $N \in \{128, 256, 512, 1024, 1536, 2048\}$ (depende ancho banda)
\item \textbf{5G NR:} $N \in \{128, 256, 512, 1024, 2048, 4096\}$ (más flexible)
\item \textbf{DVB-T (TV):} $N \in \{2048, 8192\}$ (muy robusto)
\item \textbf{DSL (ADSL2+):} $N = 512$ o $4096$
\end{itemize}

\textbf{Relación entre $N$ y robustez:}
\begin{itemize}
\item $N$ grande → subportadoras más estrechas → más robusto a delay spread largo
\item $N$ pequeño → menor latencia, menor overhead de CP
\item Trade-off: robustez vs. eficiencia/latencia
\end{itemize}

\subsection{Prefijo cíclico (CP) - Truco clave}

El prefijo cíclico es la técnica que hace funcionar OFDM en canales multitrayecto.

\textbf{¿Qué es el CP?}

Se copian las últimas $L_{CP}$ muestras del símbolo OFDM y se colocan al principio:

\textbf{Símbolo OFDM con CP:}
$$\{\underbrace{x_{N-L_{CP}}, \ldots, x_{N-1}}_{\text{CP (copia)}}, \underbrace{x_0, x_1, \ldots, x_{N-1}}_{\text{Símbolo OFDM original}}\}$$

\textbf{Longitud necesaria del CP:} $L_{CP} \geq L - 1$

donde $L$ es el número de taps del canal discreto (relacionado con el delay spread).

\textbf{¿Por qué añadir el CP?}

\textbf{Problema sin CP:} La convolución lineal del canal destruye la ortogonalidad de las subportadoras → ISI e ICI.

\textbf{Solución con CP:} El prefijo cíclico convierte la convolución lineal en convolución circular.

\textbf{Efecto mágico:} Con convolución circular, la DFT diagonaliza el canal:
\begin{itemize}
\item Cada subportadora experimenta un canal escalar $H_k$ (un solo coeficiente complejo)
\item No hay interferencia entre subportadoras
\item Ecualización trivial: $\hat{s}_k = y_k / H_k$
\end{itemize}

\textbf{Coste del CP:}
\begin{itemize}
\item Transmitimos $N + L_{CP}$ muestras pero solo $N$ contienen información nueva
\item Eficiencia: $\eta_{CP} = \frac{N}{N + L_{CP}}$
\item Ejemplo: $N=64$, $L_{CP}=16$ → $\eta_{CP} = 64/80 = 0.8$ (20\% de overhead)
\end{itemize}

\textbf{Compromiso en el diseño:}
\begin{itemize}
\item CP corto → mayor eficiencia, pero no protege contra multipaso largo
\item CP largo → protección completa, pero penaliza débito
\end{itemize}

\subsection{Parámetros}

\textbf{Espaciado subportadoras:} $\Delta f = \frac{1}{NT_s}$

\textbf{Duración símbolo OFDM:} $T_{OFDM} = NT_s$

\textbf{Duración con CP:} $T_{total} = (N + L_{CP})T_s$

\textbf{Eficiencia CP:} $\eta_{CP} = \frac{N}{N + L_{CP}}$

\textbf{Ancho de banda:} $B = N\Delta f = \frac{1}{T_s}$

\textbf{Débito total:} $D_b^{total} = \sum_{k=0}^{N-1} \frac{\log_2(M_k)}{T_{total}}$

\subsection{Ecualización en frecuencia}

$$\hat{s}_k = \frac{y_k}{H_k} = s_k + \frac{n_k}{H_k}$$

donde $H_k = \text{DFT}[h_0, \ldots, h_{L-1}, 0, \ldots, 0]_k$

Ecualización escalar (1 coeficiente complejo por subportadora).

\section{Procesos Estocásticos}

\subsection{Definiciones}

\textbf{Función expectativa:} $\mu_X(t) = E[X(t,\omega)]$

\textbf{Autocorrelación:} $R_X(t_1,t_2) = E[X(t_1,\omega)X^*(t_2,\omega)]$

\textbf{Correlación cruzada:} $R_{XY}(t_1,t_2) = E[X(t_1,\omega)Y^*(t_2,\omega)]$

\subsection{WSS (Wide-Sense Stationary)}

\begin{itemize}
\item $E[|X(t,\omega)|^2]$ finito e independiente de $t$
\item $\mu_X(t) = \mu_X$ (constante)
\item $R_X(t_1,t_2) = R_X(\tau)$ con $\tau = t_2 - t_1$
\end{itemize}

\textbf{Cicloestacionariedad:} Estadísticas periódicas con período $T_0$

\subsection{PSD y Wiener-Khinchin}

\textbf{Densidad espectral potencia:}
$$S_X(\nu) = \text{FT}[R_X(\tau)](\nu)$$

\textbf{Potencia total:}
$$P_X = R_X(0) = \int_{-\infty}^{\infty} S_X(\nu)d\nu$$

\subsection{Filtrado de procesos}

Si $Y(t) = \int h(\tau)X(t-\tau)d\tau$:

$$S_Y(\nu) = |H(\nu)|^2 S_X(\nu)$$

$$R_Y(\tau) = R_X(\tau) * h(\tau) * h^*(-\tau)$$

\section{Muestreo y Conversiones}

\subsection{Teorema de Nyquist-Shannon}

Para señal con ancho banda $B$:
$$F_s \geq 2B$$

\textbf{Reconstrucción:}
$$x(t) = \sum_{n} x_n \text{sinc}(F_s t - n)$$

con filtro pasa-bajo ideal de ancho $B \in [f_{max}, F_s - f_{max}]$

\subsection{DFT}

\textbf{Definición:}
$$\hat{z}_k = \sum_{n=0}^{N-1} z_n e^{-i2\pi nk/N}, \quad k = 0,\ldots,N-1$$

\textbf{IDFT:}
$$z_n = \frac{1}{N}\sum_{k=0}^{N-1} \hat{z}_k e^{i2\pi nk/N}, \quad n = 0,\ldots,N-1$$

\textbf{Frecuencia normalizada:} $f_k = k/N \Rightarrow \nu_k = kF_s/N$

\section{Comparación de Modulaciones}

\begin{tcolorbox}[colback=azulclaro!10,colframe=azuloscuro,title=Resumen comparativo]
\textbf{Para mismo $E_{avg}/N_0$ (alto SNR):}

M-QAM es más eficiente en potencia

Diferencia ASK vs QAM: $\approx 10\log_{10}\left(\frac{M^2-1}{2(M-1)}\right)$ dB

Diferencia ASK vs PSK: $\approx 10\log_{10}\left(\frac{M^2-1}{3\sin^2(\pi/M)}\right)$ dB

\textbf{Eficiencia espectral (mismo $\beta$):} PAM > QAM > PSK

\textbf{Recomendaciones:}
\begin{itemize}
\item QAM: mejor compromiso potencia/espectro (WiFi, LTE)
\item PSK: PAPR constante, amplificadores no lineales (satélite)
\item PAM: transmisión banda base (Ethernet)
\end{itemize}
\end{tcolorbox}

\section{Fórmulas de Referencia Rápida}

\subsection{Relaciones SNR}

\begin{align*}
E_b/N_0 &= \frac{E_{avg}}{\log_2(M)N_0}\\
E_s/N_0 &= \frac{E_{avg}}{N_0}\\
E_s/N_0 &= \log_2(M) \cdot E_b/N_0\\
\text{SNR}_{dB} &= 10\log_{10}(\text{SNR})
\end{align*}

\subsection{Logaritmos}

$$\log_2(M) = \frac{\ln(M)}{\ln(2)} = \frac{\log_{10}(M)}{\log_{10}(2)}$$

\subsection{Valores útiles Q(x)}

\begin{tabular}{ll}
$Q(0) = 0.5$ & $Q(3) \approx 1.35 \times 10^{-3}$\\
$Q(1) \approx 0.159$ & $Q(4) \approx 3.17 \times 10^{-5}$\\
$Q(2) \approx 0.023$ & $Q(5) \approx 2.87 \times 10^{-7}$
\end{tabular}

\end{multicols}

\newpage

\section{GUÍA PRÁCTICA PARA EXÁMENES}

\begin{multicols}{2}

\subsection{Metodología General}

\textbf{1. Lee TODO el problema primero}
\begin{itemize}
\item Identifica qué te dan (parámetros conocidos)
\item Identifica qué te piden (incógnitas)
\item Dibuja un esquema si es necesario
\end{itemize}

\textbf{2. Verifica unidades SIEMPRE}
\begin{itemize}
\item Frecuencias: Hz, kHz, MHz (¡convierte!)
\item Tiempos: s, ms, μs (¡cuidado con ms vs μs!)
\item Potencias: W, mW, dBm
\item Débitos: bits/s, kbits/s, Mbits/s
\end{itemize}

\textbf{3. Casos comunes en OFDM}

\textbf{Si te dan $T_{OFDM}$ y $\tau_{max}$:}
\begin{enumerate}
\item Calcula $\Delta f = 1/T_{OFDM}$
\item Calcula $T_{CP} \geq \tau_{max}$ (típicamente $T_{CP} = \tau_{max}$)
\item Calcula $T_{total} = T_{OFDM} + T_{CP}$
\item Para débito, ¡USA $T_{total}$!
\end{enumerate}

\textbf{Si preguntan por pilotos:}
\begin{enumerate}
\item \textbf{Frecuencia:} $\Delta f_{pilot} \leq \frac{1}{2\tau_{max}}$
\item Calcula: $N_{max} = \lfloor \frac{1}{2\tau_{max} \cdot \Delta f} \rfloor$
\item Elige la opción que cumpla $N \leq N_{max}$
\item Si hay empate, elige la que da mejor débito
\end{enumerate}

\begin{enumerate}[resume]
\item \textbf{Tiempo:} $T_{pilot} \leq \frac{T_{coh}}{2}$
\item Calcula: $N_{OFDM,max} = \lfloor \frac{T_{coh}}{2T_{OFDM}} \rfloor$
\item Si $T_{OFDM} \approx T_{coh}$: pilotos en TODOS los símbolos
\end{enumerate}

\textbf{4. Ecualizadores - Qué responder}

\textbf{Pregunta DFE:}
\begin{itemize}
\item Esquema: Feedforward + Feedback + Decisor
\item Ventaja: No amplifica ruido en feedback
\item Desventaja: Propagación de errores
\item Solución: Training sequences + código corrector
\end{itemize}

\textbf{Pregunta Viterbi/ML:}
\begin{itemize}
\item Dibuja trellis con estados = memoria canal
\item Para $y_n = Aa_n + Aa_{n-1}$: 2 estados $\{\pm 1\}$
\item Menciona: Métrica de rama, sobrevivientes, traceback
\end{itemize}

\subsection{Ejercicios Tipo por Tema}

\subsubsection{Tipo 1: Canal Discreto Equivalente}

\textbf{Datos típicos:}
\begin{itemize}
\item Filtro emisión: $h_e(t) = \mathds{1}_{[0,T]}$
\item Canal: $c(t) = \delta(t-\tau_1) + \alpha_2\delta(t-\tau_2) + \ldots$
\item Matched filter: $g_r(t) = h_e(-t)$
\end{itemize}

\textbf{Pasos:}
\begin{enumerate}
\item Calcula $r(t) = (h_e * g_r)(t)$ → Triangular
\item Calcula $(r * c)(t)$ → Suma de triángulos desplazados
\item Muestrea en $t_0 + nT$ → Coeficientes $h_n$
\item Identifica ISI: pre-cursor y post-cursor
\item Propón ecualizador según ISI
\end{enumerate}

\textbf{Fórmula clave para matched filter:}
$$r(t) = \Lambda\left(\frac{t-T}{T}\right) \cdot T = \begin{cases}
t & 0 \leq t < T\\
2T-t & T \leq t < 2T\\
0 & \text{otro caso}
\end{cases}$$

\subsubsection{Tipo 2: OFDM Submarino / Móvil}

\textbf{Datos típicos:}
\begin{itemize}
\item $N$ subportadoras, PSK-$M$, código $r_c$
\item $\tau_{max}$, $T_{coh}$, $T_{OFDM}$
\item Pilotos: elegir configuración
\end{itemize}

\textbf{Pasos sistemáticos:}
\begin{enumerate}
\item $\Delta f = 1/T_{OFDM}$ (ortogonalidad)
\item $T_{CP} = \tau_{max}$ (combatir ISI)
\item $T_{total} = T_{OFDM} + T_{CP}$ (tiempo real)
\item Pilotos frecuencia: $N_{pilot,f} \leq 1/(2\tau_{max}\Delta f)$
\item Pilotos tiempo: Comparar $T_{OFDM}$ con $T_{coh}/2$
\item Bits/OFDM: $(N - N_{pilotos}) \times \log_2(M) \times r_c$
\item Débito: Bits / $T_{total}$ (¡no $T_{OFDM}$!)
\end{enumerate}

\textbf{Trampa común:} Usar $T_{OFDM}$ en vez de $T_{total}$ para débito.

\textbf{Eficiencia CP:}
$$\eta_{CP} = \frac{T_{OFDM}}{T_{total}} = \frac{T_{OFDM}}{T_{OFDM}+T_{CP}}$$

Overhead: $1 - \eta_{CP} = T_{CP}/T_{total}$

\subsubsection{Tipo 3: Estimación de Canal}

\textbf{Preguntas típicas:}
\begin{itemize}
\item Calcular intercorrelación $R_{ya}(u)$
\item Proponer método de estimación
\item Longitud óptima de training
\end{itemize}

\textbf{Fórmulas esenciales:}
$$R_{ya}(u) = \sum_k y_k a^*(k-u) = h * R_{aa}$$

Si $R_{aa}(u) = \delta(u)$ (secuencia ideal):
$$\hat{h}_u = R_{ya}(u)$$

\textbf{Solución matricial (LS):}
$$\mathbf{y} = \mathbf{A}\mathbf{h} + \mathbf{b}$$
$$\hat{\mathbf{h}} = (\mathbf{A}^H\mathbf{A})^{-1}\mathbf{A}^H\mathbf{y}$$

\textbf{Training óptimo:}
\begin{itemize}
\item Longitud $\geq$ memoria canal
\item Balancear: overhead vs calidad estimación
\item Repetir cada $T_{coh}/2$ si canal varía
\end{itemize}

\subsubsection{Tipo 4: Códigos Convolucionales}

\textbf{Datos típicos:}
\begin{itemize}
\item Registros: $K$ (típicamente 3-7)
\item Rendimiento: $r_c = k/n$ (ejemplo: 1/2, 2/3)
\item Generadores: $(g_1, g_2, \ldots)$ en octal
\end{itemize}

\textbf{Pasos codificación:}
\begin{enumerate}
\item Inicializa registros a 0
\item Por cada bit entrada $b_k$:
   \begin{itemize}
   \item Desplaza registros
   \item Calcula salidas según generadores
   \item Guarda bits codificados
   \end{itemize}
\item Flush: añade 0s para vaciar registros
\end{enumerate}

\textbf{Ejemplo $(7,5)_8$ con $r_c=1/2$:}
\begin{itemize}
\item $7_8 = 111_2$: $c_1 = b_k \oplus b_{k-1} \oplus b_{k-2}$
\item $5_8 = 101_2$: $c_2 = b_k \oplus b_{k-2}$
\end{itemize}

\textbf{Matriz de paridad:}

Para verificar ortogonalidad: $\mathbf{m} \cdot \mathbf{H}^T = \mathbf{0}$

\subsubsection{Tipo 5: Probabilidad de Error con ISI}

\textbf{Modelo típico:}
$$y_n = Aa_n + Aa_{n-1} + b_n$$

con $a_n \in \{-1, +1\}$ y decisión en 0.

\textbf{Análisis sistemático:}
\begin{enumerate}
\item Lista todos los casos: $(a_n, a_{n-1})$
\item Calcula $y_n$ para cada caso (sin ruido)
\item Determina decisión según umbral
\item Identifica casos de error
\item Calcula $P_e$ asumiendo equiprobables
\end{enumerate}

\textbf{Resultado típico:}
$$P_e = \frac{1}{4} \quad \text{(error cuando } a_n=-1, a_{n-1}=+1)$$

\textbf{Concepto clave:} Incluso sin ruido ($\sigma^2 \to 0$), ISI causa errores irreductibles.

\subsection{Checklist Final Antes del Examen}

\begin{tcolorbox}[colback=verdeclaro!10,colframe=azuloscuro,title=Verifica que sabes]

\textbf{OFDM:}
\begin{itemize}
\item[$\square$] Diferencia $T_{OFDM}$ vs $T_{total}$
\item[$\square$] Criterios Nyquist (frecuencia y tiempo)
\item[$\square$] Cálculo débito con CP y código
\item[$\square$] Configuración óptima pilotos
\end{itemize}

\textbf{Ecualizadores:}
\begin{itemize}
\item[$\square$] Dibujar DFE completo
\item[$\square$] Ventajas/desventajas DFE vs ML
\item[$\square$] Trellis para Viterbi (2 estados básico)
\item[$\square$] Soluciones propagación errores
\end{itemize}

\textbf{Canal:}
\begin{itemize}
\item[$\square$] Convolución triangular (matched filter)
\item[$\square$] Muestreo y coeficientes $h_n$
\item[$\square$] Delay spread $\tau_{max} - \tau_{min}$
\item[$\square$] ISI pre-cursor vs post-cursor
\end{itemize}

\textbf{Estimación:}
\begin{itemize}
\item[$\square$] $R_{ya} = h * R_{aa}$
\item[$\square$] Solución LS matricial
\item[$\square$] Training óptimo vs overhead
\end{itemize}

\textbf{Códigos:}
\begin{itemize}
\item[$\square$] Codificador convolucional (diagrama)
\item[$\square$] Octal → binario → XOR
\item[$\square$] Rendimiento $r_c = k/n$
\item[$\square$] Matriz paridad $\mathbf{H}$
\end{itemize}

\end{tcolorbox}

\subsection{Errores Comunes a Evitar}

\begin{tcolorbox}[colback=rojoclaro!10,colframe=rojoclaro!70!black,title=¡CUIDADO!]

\textbf{Error 1: Débito OFDM}
\begin{itemize}
\item ✗ Usar $R_b = \text{bits}/T_{OFDM}$
\item ✓ Usar $R_b = \text{bits}/T_{total}$ con $T_{total} = T_{OFDM} + T_{CP}$
\end{itemize}

\textbf{Error 2: Pilotos límite}
\begin{itemize}
\item ✗ Decir "no cumple" cuando $T_{OFDM} = T_{coh}$
\item ✓ Opción (a) es "caso límite aceptable"
\end{itemize}

\textbf{Error 3: Código corrector}
\begin{itemize}
\item ✗ Multiplicar por $n/k$ (al revés)
\item ✓ Multiplicar por $r_c = k/n$ para bits info
\end{itemize}

\textbf{Error 4: Unidades tiempo}
\begin{itemize}
\item ✗ Mezclar ms con μs sin convertir
\item ✓ Todo a segundos antes de calcular
\end{itemize}

\textbf{Error 5: Matched filter}
\begin{itemize}
\item ✗ Pensar que $r(t)$ es rectangular
\item ✓ $r(t)$ es TRIANGULAR (convolución de 2 rectangulares)
\end{itemize}

\textbf{Error 6: ISI y ruido}
\begin{itemize}
\item ✗ Pensar que sin ruido $P_e = 0$
\item ✓ ISI causa errores incluso sin ruido
\end{itemize}

\end{tcolorbox}

\subsection{Fórmulas para Copiar en Examen}

\textbf{OFDM básico:}
\begin{align*}
\Delta f &= \frac{1}{T_{OFDM}} & T_{total} &= T_{OFDM} + T_{CP}\\
T_{CP} &\geq \tau_{max} & R_b &= \frac{\text{bits info}}{T_{total}}
\end{align*}

\textbf{Pilotos:}
\begin{align*}
\Delta f_{pilot} &\leq \frac{1}{2\tau_{max}} & T_{pilot} &\leq \frac{T_{coh}}{2}
\end{align*}

\textbf{Bits OFDM:}
\begin{align*}
N_{bits} &= N_{data} \times \log_2(M) \times r_c\\
N_{data} &= N_{total} - N_{pilotos}
\end{align*}

\textbf{Canal:}
\begin{align*}
y_n &= \sum_{l=0}^{L-1} h_l a_{n-l} + b_n\\
R_{ya}(u) &= h * R_{aa}
\end{align*}

\textbf{Matched filter rectangular:}
$$r(t) = \begin{cases}
t & 0 \leq t < T\\
2T-t & T \leq t < 2T\\
0 & \text{otro}
\end{cases}$$

\end{multicols}

\end{document}
